\chapter{Unbiased updates}
\label{ch: unbiased}

% The previous chapter established that 
% MC-O-PI converges to the optimal solution if the sequences $(q_k)$ and $(\pol_k)$ generated by exact MC-O-PI converge to $\qopt$ and $\best{\pol}$.
% Crucially, the parameters controlling these sequences are the update distributions $(\update_k)$.


The previous chapter established that MC-O-PI solutions remain trapped in a shrinking tube around their reinitialised mean-field dynamics.
Building on this insight, we extend
the results of \citet{Tsitsiklis2002OnIteration}, \citet{Chen2018OnProblem} and \citet{Liu2021OnStarts} to a broader and more practical class of update distributions.
Specifically, while they guarantee convergence when updates are \textit{uniform over all state-action pairs} (\autoref{asp: unbiased tsitsiklis}), we show that convergence still holds when updates are only \textit{uniform over all actions in each state}.

\begin{definition}[Uniform updates over actions in each state]
\label{asp: unbiased over policies}
The update distribution $\update$ is uniform over all actions in each state,
% The update map $\updateset$ yields update distributions $\update$ that are uniform over the policies, i.e.
if
for every state $s \in \sset$ and all actions $a, a' \in \aset(s)$ we have
\begin{align}
    % \label{eq: uniform bias over policies}
    \nonumber
    % \forall k \geq 0,
    % \forall \pol \in \polset,
    % \forall s \in \sset,
    % \forall a, a' \in \aset(s)
    % : \quad
    \update(s, a) = \update(s, a') .
\end{align}
% MC-O-PI thus updates the actions of each state as frequently.
\end{definition}

\vspace{-0.6cm}

\autoref{asp: unbiased over policies} generalises \autoref{asp: unbiased tsitsiklis} because it allows updating different states at different frequencies.
For instance, under initial-visit updates, the initial state of each episode can now be sampled according to \textit{any} distribution provided the initial action in that state is sampled uniformly.
This additional flexibility disrupts Tsitsiklis' proof as the update distributions no longer commute with the transition operator in \eqref{eq: tsitsiklis commutation}. However, we show that MC-O-PI still converges to the optimal solution.

This weaker condition is more practical for real-world applications, in particular when the state-space is large.
For instance, chess has $O(10^{45})$ different board layouts but only $O (10^1)$ legal moves per layout.
Uniformly sampling over the joint state-action space is thus computationally intractable, but sampling an action uniformly from a chosen state is trivial.
Henceforth, \textit{uniform updates} refer to \autoref{asp: unbiased over policies} rather than \autoref{asp: unbiased tsitsiklis}.


This chapter
first demonstrates that the mean-field dynamics under uniform updates generate policies that improve monotonically.
We then show that the stochastic perturbations in MC-O-PI cannot consistently prevent this improvement, and hence that MC-O-PI converges to the optimal action values.
Since the non-contracting, discontinuous and non-convex MC-O-PI dynamics preclude using classical methods (see \autoref{sec: limitations of existing methods}),
% we turn the mean-field improvement into a new probabilistic lock-in argument from first principles.
we derive convergence from first principles.

% As described in \autoref{sec: track mean field},
% studying the asymptotic behaviour of MC-O-PI comes down to understanding how the target action values change over time.
% To isolate these changes, we introduce the transition times.


% We first show that if the update distributions are uniform,
% MC-O-PI transitions, on average, to monotonically improving policies.
% This insight allows us to extend previous convergence results while clarifying the convergence mechanism.


% %%%%%%%%%%%%%%%%%%%%%%%%%%%%%%%%%%%%%%%%%
% %%%%%%%%%%%%%%%%%%%%%%%%%%%%%%%%%%%%%%%%%
% \section{Experiment}
% \label{sec: uniform experiment}

% We show that the additional flexibility that \autoref{asp: unbiased tsitsiklis} offers over \autoref{asp: unbiased over policies} does not affect the qualitative behaviour of MC-O-PI.
% Namely, both move along straight lines towards the current action values.

% % show experiment with synchronous and asynchronous updates

% %%%%%%%%%%%%%%%%%%%%%%%%%%%%%%%%%%%%%%%%%
% \subsection{MDP}

% Consider the discounted, two-state, deterministic Markov Decision Process depicted in \autoref{fig: experiment uniform bias MDP}. In this MDP, each state offers two possible actions: ``stay,'' which yields a reward of 0, and ``move,'' which yields a reward of 1. The optimal policy in this scenario is to always choose the ``move'' action and transition to the other state, while the worst policy is to always choose the ``stay'' action. For clarity, in each state we denote the optimal action as $\best{a}$ and the suboptimal action as $\worst{a}$.

% \begin{figure}[htbp]
    \centering
    \begin{subfigure}[b]{0.48\textwidth}
        \centering
        \begin{tikzpicture}[
            scale=1, auto, node distance=2cm, 
            every edge/.style={draw, <-, font=\small},
            state/.style={circle, fill=mylightgrey, draw=black, text=black}
        ]
            % Nodes
            \node[state] (s1) at (0,0) {$s_1$};
            \node[state] (s2) at (2,0) {$s_2$};
            % Edges
            \path[->]
                (s1) edge[out=160, in=200, loop, looseness=8] node[pos=0.5, left] {0} (s1)
                (s2) edge[bend right=20] node[pos=0.5, above] {1} (s1)

                (s2) edge[out=20, in=340, loop, looseness=8] node[pos=0.5, right] {0} (s2)
                (s1) edge[bend right=20] node[pos=0.5, below] {1} (s2);
        \end{tikzpicture}
        \caption{MDP}
    \end{subfigure}%
   
    \vfill
    \vspace{1cm}

    \begin{subfigure}[b]{0.48\textwidth}
        \centering
        \begin{tikzpicture}[
            scale=1, auto, node distance=2cm, 
            every edge/.style={draw, <-, font=\small},
            state/.style={circle, fill=mylightgrey, draw=black, text=black}
        ]
            % Nodes
            \node[state] (s1) at (0,0) {$s_1$};
            \node[state] (s2) at (2,0) {$s_2$};
            % Edges
            \path[->]
                (s1) edge[out=160, in=200, loop, looseness=8, draw=mylightgrey] node[pos=0.5, left, text=mylightgrey] {0} (s1)
                (s2) edge[bend right=20] node[pos=0.5, above] {1} (s1)

                (s2) edge[out=20, in=340, loop, looseness=8, draw=mylightgrey] node[pos=0.5, right, text=mylightgrey] {0} (s2)
                (s1) edge[bend right=20] node[pos=0.5, below] {1} (s2);
        \end{tikzpicture}
        \caption{Policy $\best{\pol}$}
    \end{subfigure}%
    \hfill
    \begin{subfigure}[b]{0.48\textwidth}
        \centering
        \begin{tikzpicture}[
            scale=1, auto, node distance=2cm, 
            every edge/.style={draw, <-, font=\small},
            state/.style={circle, fill=mylightgrey, draw=black, text=black}
        ]
            % Nodes
            \node[state] (s1) at (0,0) {$s_1$};
            \node[state] (s2) at (2,0) {$s_2$};
            % Edges
            \path[->]
                (s1) edge[out=160, in=200, loop, looseness=8] node[pos=0.5, left] {0} (s1)
                (s2) edge[bend right=20, draw=mylightgrey] node[pos=0.5, above, text=mylightgrey] {1} (s1)

                (s2) edge[out=20, in=340, loop, looseness=8] node[pos=0.5, right] {0} (s2)
                (s1) edge[bend right=20, draw=mylightgrey] node[pos=0.5, below, text=mylightgrey] {1} (s2);
        \end{tikzpicture}
        \caption{Policy $\worst{\pol}$}
    \end{subfigure}%
   
    \vfill
    \vspace{1cm}

    \begin{subfigure}[b]{0.48\textwidth}
        \centering
        \begin{tikzpicture}[
            scale=1, auto, node distance=2cm, 
            every edge/.style={draw, <-, font=\small},
            state/.style={circle, fill=mylightgrey, draw=black, text=black}
        ]
            % Nodes
            \node[state] (s1) at (0,0) {$s_1$};
            \node[state] (s2) at (2,0) {$s_2$};
            % Edges
            \path[->]
                (s1) edge[out=160, in=200, loop, looseness=8, draw=mylightgrey] node[pos=0.5, left, text=mylightgrey] {0} (s1)
                (s2) edge[bend right=20] node[pos=0.5, above] {1} (s1)

                (s2) edge[out=20, in=340, loop, looseness=8] node[pos=0.5, right] {0} (s2)
                (s1) edge[bend right=20, draw=mylightgrey] node[pos=0.5, below, text=mylightgrey] {1} (s2);
        \end{tikzpicture}
        \caption{Policy $\pol_1$}
    \end{subfigure}%
    \hfill
    \begin{subfigure}[b]{0.48\textwidth}
        \centering
        \begin{tikzpicture}[
            scale=1, auto, node distance=2cm, 
            every edge/.style={draw, <-, font=\small},
            state/.style={circle, fill=mylightgrey, draw=black, text=black}
        ]
            % Nodes
            \node[state] (s1) at (0,0) {$s_1$};
            \node[state] (s2) at (2,0) {$s_2$};
            % Edges
            \path[->]
                (s1) edge[out=160, in=200, loop, looseness=8] node[pos=0.5, left] {0} (s1)
                (s2) edge[bend right=20, draw=mylightgrey] node[pos=0.5, above, text=mylightgrey] {1} (s1)

                (s2) edge[out=20, in=340, loop, looseness=8, draw=mylightgrey] node[pos=0.5, right, text=mylightgrey] {0} (s2)
                (s1) edge[bend right=20] node[pos=0.5, below] {1} (s2);
        \end{tikzpicture}
        \caption{Policy $\pol_2$}
    \end{subfigure}
    \caption{MDP description}
    \label{fig: experiment uniform bias MDP}
\end{figure}

% %%%%%%%%%%%%%%%%%%%%%%%%%%%%%%%%%%%%%%%%%
% \subsection{MC-O-PI}

% explain that update one state twice as often as the other = satisfies \autoref{asp: unbiased over policies} but not \autoref{asp: unbiased tsitsiklis}

% % simulation with uniform bias over state-actions

% % simulation with uniform bias over policies


% \begin{figure}
%     \centering
%     \begin{subfigure}{0.48\textwidth}
%         \centering
%         % \includegraphics[]{}
%         \caption{simulation with uniform bias over state-actions}
%         \label{fig: uniform simulation over all state-actions}
%     \end{subfigure}
%     \hfill
%     \begin{subfigure}{0.48\textwidth}
%         \centering
%         % \includegraphics[]{}
%         \caption{simulation with uniform bias over policies}
%         \label{fig: uniform simulation over policies}
%     \end{subfigure}
%     \caption{no qualitative difference}
%     \label{fig: uniform simulation}
% \end{figure}
% %%%%%%%%%%%%%%%%%%%%%%%%%%%%%%%%%%%%%%%%%
% \subsection{Discussion}

% \textcolor{red}{different bias but similar behaviour => same underlying phenomenon?} \\
% - \autoref{fig: uniform simulation} shows that even though MC-O-PI with uniform bias over the policies does not satisfy the conditions of \autoref{thm: tsitsiklis convergence proof}, its behaviour is very similar to the behaviour when the update is uniform over all state-action pairs
% - Namely, solution moves along straight lines towards the action-values and does not present sliding mode around the boundary.
% - Similar behaviour suggests the same underlying phenomenon.


% The exploring starts condition ensures that the resulting update distributions are strictly positive for all state-action pairs.

% are essentially the same as the modified robbins-Monro conditions in \autoref{asp: stochastic conditions on effective learning rate}
% for clarity we decided to separate them in the following discussion


% Note that the policy itself need not be constant on such a block.

% In other words, the definition of the transition times does not mean that the greedy policy is unchanged between times $t_n$ and $t_{n+1}$, but rather that the target action values remain unchanged, and thus also the behaviour of MC-O-PI.
% - noise and update distribution might be affected by policy



% Suppose that \eqref{eq: strict improvement theorem} were an equality for every state. The Policy Improvement Theorem would then give $v_{\pi} = v_{\pi'}$, and consequently $q_{\pi} = q_{\pi'}$, since for every state $s \in \mathcal S$ and action $a \in \mathcal A(s)$,
% \[
% q_{\pi}(s,a) - q_{\pi'}(s,a) = \gamma \sum_{s' \in \mathcal S} p(s' \mid s,a) \big( v_{\pi}(s') - v_{\pi'}(s') \big) .
% \]
% However, this contradicts the assumption that $q_\pi \neq q_{\pi'}$.
% Therefore, \eqref{eq: strict improvement theorem} must be strict in at least one state, and hence $\pi \prec \pi'$.

%%%%%%%%%%%%%%%%%%%%%%%%%%%%%%%%%%%%%%%%%
%%%%%%%%%%%%%%%%%%%%%%%%%%%%%%%%%%%%%%%%%
\section{Mean-field generates monotonically improving policies}

Given that the stochastic perturbations in initial- or first-visit MC-O-PI are unbiased (\autoref{thm: properties of noise}),
the mean-field update that arises when averaging out these perturbations is
\begin{align}
\label{eq: deterministic mcopi}
    q_{k+1}
    \in
    \big\{ q_k + \lrate_k \update_k ( q_{\pi_k} - q_k )
    :
    \pol_k \in \grpol(q_k),
    \update_k = \update(\start_k, \pol_k, f_k)
    \big\}
\end{align}
with learning rates $(\lrate_k)_{k\ge 0}$, start distributions $(\start_k)_{k\ge 0}$ and update functions $(f_k)_{k\ge 0}$.

As established in the previous chapter, the asymptotic behaviour of this system hinges on the evolution of the target action values $q_{\pi_k}$.
Using the transition times $(t_n)_{n \ge 0}$ from \autoref{def: transition times}, we show that
if the update distributions $(\update_k)_{k\ge 0}$ remain uniform over all actions in each state,
the sequence $(q_{\pi_{t_n}})_{n \ge 0}$ increases monotonically up to $\qopt$,
and hence the estimate $q_k$ converges to the optimal action values.

We rely on the following result at the transition times.

\begin{lemma}
\label{thm: strict improvement theorem}
For all policies $\pi, \pi' \in \polset$ that satisfy
\begin{align}
\label{eq: strict improvement theorem}
q_\pi(s,\pi) \le q_\pi(s,\pi')
\end{align}
for every state $s \in \sset$,
if $q_\pi \neq q_{\pi'}$, then policy $\pi'$ is better than $\pi$ because $v_{\pi} \le v_{\pi'}$ and $v_{\pi}(s) < v_{\pi'}(s)$ for at least one state $s \in \sset$.
\end{lemma}
\vspace{-0.7cm}
\begin{proof}
By the \hyperref[thm: policy improvement theorem]{Policy Improvement Theorem}, inequality \eqref{eq: strict improvement theorem} implies that $v_\pi \le v_{\pi'}$.
% If $v_{\pi} = v_{\pi'}$, then for every $s \in \mathcal S$ and $a \in \mathcal A(s)$,
Suppose for contradiction that $v_{\pi} = v_{\pi'}$. Then for every state $s \in \mathcal S$ and action $a \in \mathcal A(s)$,
\begin{align*}
q_{\pi}(s,a) 
&= r(s,a) + \gamma \sum_{s_+ \in \mathcal S} \transit (s_+ \mid s,a) v_{\pi}(s_+)
\\
&= r(s,a) + \gamma \sum_{s_+ \in \mathcal S} \transit (s_+ \mid s,a) v_{\pi'}(s_+)
\\
&= q_{\pi'}(s,a).
\end{align*}
This contradicts the assumption that $q_\pi \neq q_{\pi'}$. Therefore, there exists at least one state $s \in \mathcal S$ where $v_{\pi}(s) < v_{\pi'}(s)$, implying that policy $\pi'$ is better than $\pi$.
\end{proof}


% This section first formalises the concept of \textit{policy transition},
% then proves the strict improvement at each transition,
% and finally establishes the convergence of \eqref{eq: deterministic mcopi}.


%%%%%%%%%%%%%%%%%%%%%%%%%%%%%%%%%%%%%%%%%
\subsection{Policy improvement at every transition}


% For every $q \in \qset$, a mean-field update of MC-O-PI is
% choose a greedy policy $\pol \in \grpol(q)$, $\update = \update(\start, \pol, f)$
% update as
% \begin{align}
% \label{eq: definition of q+ mean field uniform section}
%     q_+ 
%     =
%     q + \lrate \update (\qpol - q)
% \end{align}

By showing that inequality \eqref{eq: strict improvement theorem} holds at every transition time, we now prove that policies generated by the mean-field dynamics in \eqref{eq: deterministic mcopi} improve monotonically.


\begin{proposition}
\label{thm: unbiased mces policy improves}
Let \((\pi_k)_{k\ge 0}\) be a sequence of policies generated by the mean-field dynamics in \eqref{eq: deterministic mcopi}
with
sampling distributions $(\sigma_k)_{k \ge 0}$ that satisfy exploring starts (\autoref{def: exploring starts}),
and update distributions $(\mu_k)_{k \ge 0}$ that are uniform over all actions in each state (\autoref{asp: unbiased over policies}).
Let \((t_n)_{n\ge 0}\) be the associated transition times (\autoref{def: transition times}).
% Under \autoref{asp: standing},
Then, for every \(n \in \mathbb Z_{\ge 0}\) with \(t_{n+1}<\infty\), policy
\(\pi_{t_{n+1}}\) is better than \(\pi_{t_n}\).

% Given a point $q \in \qset$, choose any policy $\pi \in \grpol(q)$ and define
% \begin{align}
% \label{eq: definition of q+ uniform section}
%     q_+
%     \coloneqq
%     % \in
%     % \big\{ 
%     q + \lrate \update (\qpol - q) .
%     % :
%     % \pol \in \grpol(q),
%     % \update = \update(\start, \pol, f)
%     % \big\} .
% \end{align}
% For every $q \in \qset$,
% Let $q_{k+1}$ satisfy \eqref{eq: deterministic mcopi} with policy $\pi_k$.

% \begin{align}
% % \label{eq: pit in unbiased lemma for improvement}
%     q_{\pol_k}(s, {\pol_k}) \le q_{\pol_k}(s, \pol_{k+1})
%     % \qquad
%     % \forall \, s \in \sset .
% \end{align}
% for every state $s \in \mathcal S$.

% Suppose $q$ and $q_+$ satisfy \eqref{eq: definition of q+ uniform section}.
% If the update distribution $\update$ is uniform over all actions in each state (\autoref{asp: unbiased over policies}),
% % (\autoref{asp: unbiased over policies}),
% then for all policies $\pol \in \grpol(q)$ and $\pol_+ \in \grpol(q_+)$, it follows that $\pol \preceq \pol_+$.
% Then, for all policies $\pol \in \grpol(q)$ and $\pol_+ \in \grpol(q_+)$, it holds that
% \begin{align}
%     v_{\pol} \leq v_{\pol_+} .
% \end{align}
% In other words, the policy $\pol_+$ is as good as, or better than, $\pol$.
% policy $\pol_+$ is at least as good as policy $\pol$.
\end{proposition}
\vspace{-0.6cm}
\begin{proof}
Let $(q_k)_{k \ge 0}$ be a solution of \eqref{eq: deterministic mcopi} generated with update distributions $(\mu_k)_{k \ge 0}$, and let $(\pi_k)_{k \ge 0}$ be the resulting policy sequence with associated transition times $(t_n)_{n \ge 0}$.
\autoref{thm: all inequalities when switching policy}(1) shows that for every time $k \ge 0$ and every state $s \in \sset$,
\begin{multline}
\label{eq: inequality at transition}
\update_k(s, \pol_k) \big( q_{\pol_k}(s, \pol_k) - q_k(s, \pol_k) \big)
\\
    \leq
    \update_k(s, \pol_{k+1}) \big( q_{\pol_k}(s, \pol_{k+1}) - q_k(s, \pol_k) \big) .
\end{multline}
By assumption, the update distributions are strictly positive and uniform over all actions in each state, meaning
\(\update_k(s,\pi_k)=\update_k(s,\pi_{k+1})>0\).
Inequality \eqref{eq: inequality at transition} thus reduces to
\begin{align}
    \label{eq: ineq for improvement}
    q_{\pi_k}(s,\pi_k) \le q_{\pi_k}(s,\pi_{k+1}).
\end{align}
%Since $\update_k(s,\cdot)$ is constant and strictly positive on $\mathcal A(s)$,
%rearranging the previous inequality yields
%\begin{align*}
   % q_{\pol_k}(s, \pol_k) \le q_{\pol_k}(s, \pol_{k+1}) %.
%\end{align*}

Now choose an $n \ge 0$ such that $t_{n+1}<\infty$. Applying inequality \eqref{eq: ineq for improvement} at time $k=t_{n+1}-1$ and using the fact that $q_{\pi_{t_{n+1}-1}}=q_{\pi_{t_n}}$, we obtain
\begin{align*}
    q_{\pi_{t_{n}}}(s,\pi_{t_n}) \le q_{\pi_{t_n}}(s,\pi_{t_{n+1}}).
\end{align*}
By definition of the transition times, we also know that $q_{\pi_{t_n}} \neq q_{\pi_{t_{n+1}}}$.
As a result, \autoref{thm: strict improvement theorem} implies that policy $\pi_{t_{n+1}}$ is better than $\pi_{t_n}$.
\end{proof}
% Since this inequality holds for every state $s \in \sset$,
% the \hyperref[thm: policy improvement theorem]{Policy Improvement Theorem} indicates that $\pol \preceq \pol_+$.
% \autoref{thm: all inequalities when switching policy}(1) shows that 
% for every policy 
% % $\pol \in \grpol(q)$, 
% $\pol_+ \in \grpol(q_+)$
% and state $s \in \sset$ 
% we have
% \begin{align}
% \label{eq: inequality at transition on pol}
%     \update(s, \pol) ( \qval_{\pol}(s, \pol) - \qval(s, \pol) )
%     \leq
%     \update(s, \pol_+) ( \qval_{\pol}(s, \pol_+) - \qval(s, \pol) ) .
% \end{align}
% Since the update distribution $\update$ is uniform over all actions in each state
% % \autoref{asp: unbiased over policies} implies that
% % the biases towards policies $\pol$ and $\pol_+$ are equal, namely 
% % $\update(s, \pol) = \update(s, \pol_+)$.
% inequality \eqref{eq: inequality at transition on pol} simplifies to
% \begin{align}
% \label{eq: pit in unbiaed lemma for improvement}
%     \qval_\pol(s, \pol) \leq \qval_\pol(s, \pol_+)
% \end{align}
% because $\update(s, \pol) = \update(s, \pol_+)$ for every state.
% % As a result, 
% % since inequality \eqref{eq: pit in unbiaed lemma for improvement} holds for every state, 
% As a result, the \hyperref[thm: policy improvement theorem]{Policy Improvement Theorem} indicates that $\pol \preceq \pol_+$.
% \end{proof}

\autoref{thm: unbiased mces policy improves} indicates that 
MC-O-PI with uniform updates would visit improving policies
if stochastic perturbations were averaged out, for instance by approximating $q_{\pi_k}$ as the average return from many episodes.
% 
This conclusion remains valid when policies are stochastic, as long as for every state $s \in \sset$,
\begin{align*}
    % \begin{cases}
        &\sum_{a \in \aset(s)} \pol_k(a \mid s) q_k(s,a) 
        \geq 
        \sum_{a \in \aset(s)} \pol_{k+1}(a \mid s) q_k(s,a) ,
        \\
        &\sum_{a \in \aset(s)} \pol_k(a \mid s) q_{k+1}(s,a) 
        \leq 
        \sum_{a \in \aset(s)} \pol_{k+1}(a \mid s) q_{k+1}(s,a) ,
    % \end{cases}    
\end{align*}
such as when the policies $(\pol_k)$ are epsilon-greedy on the estimates $(q_k)$.

% Technically, this result does not require exploring starts. It suffices to impose Robbins-Monro learning rates and ensure that each state-action pair is updated infinitely often. We discuss this relaxation in section \ref{sec: relaxing conditions}.
% 

\begin{remark}
Tsitsiklis' result agrees with \autoref{thm: unbiased mces policy improves} when update distributions are uniform over all state-action pairs.
Namely, the key insight of the proof is that for all $q_k$ and $q_{k+1}$ that satisfy \eqref{eq: deterministic mcopi},
\begin{align*}
    q_{k+1} - \B q_{k+1} \leq (\I - \lrate_k \update_k) (q_k - \B q_k) 
\end{align*}
where $\B$ is the Bellman operator from \autoref{def: bellman operators}.
Since the effective learning rate $\lrate_k \update_k$ is eventually smaller than one for all state-action pairs, 
we obtain
for every policy $\pol_{k+1} \in \grpol(q_{k+1})$,
\begin{align*}
    &q_{k+1} - \B_{\pol_{k+1}} q_{k+1}
    = q_{k+1} - \B q_{k+1}
    \\
    & \hspace{1.7cm} \leq (\I - \lrate_k \update_k) (q_k - \B q_k)
    \leq (\I - \lrate_k \update_k) (q_k - \B_{\pol_{k+1}} q_k ) .
\end{align*}
After rearranging the inequality, expanding $q_{k+1}$ as in \eqref{eq: deterministic mcopi},
and using the affine properties of the Bellman operator,
we observe that
\begin{align*}
    0 &\geq q_{k+1} - \B_{\pol_{k+1}} q_{k+1} - (\I - \lrate_k \update_k) ( q_k - \B_{\pol_{k+1}} q_k )
    \\
    &= (\I-\lrate_k \update_k) q_k + \lrate_k \update_k q_{\pol_k}
    - \B_{\pol_{k+1}} \big( (\I-\lrate_k \update_k) q_k + \lrate_k \update_k q_{\pol_k} \big)
    \\
    &\quad - (\I-\lrate_k \update_k) ( q_k - \B_{\pol_{k+1}} q_k)
    \\
    &= \lrate_k \update_k q_{\pol_k} - (\I-\lrate_k \update_k) \B_{\pol_{k+1}} q_k 
    \\
    &\quad - \lrate_k \update_k \B_{\pol_{k+1}}  q_{\pol_k} + (\I-\lrate_k \update_k) \B_{\pol_{k+1}} q_k
    \\
    % &= \lrate \update q_{\pol} - (\I-\lrate \update) \B_{\pol_+} q - \lrate \update \B_{\pol_+} q_{\pol} + (\I-\lrate \update) \B_{\pol_+} q
    % \\
    &= \lrate_k \update_k ( q_{\pol_k} - \B_{\pol_{k+1}} q_{\pol_k} )
\end{align*}
Since $\lrate_k \update_k$ is always positive under exploring starts, we see that $q_{\pol_k} \leq \B_{\pol_{k+1}} q_{\pol_k}$.
The monotone and contracting properties of the Bellman operator then imply that 
\begin{align*}
    q_{\pol_k} 
    \leq \B_{\pol_{k+1}} q_{\pol_k} 
    \leq \B_{\pol_{k+1}}^2 q_{\pol_k}
    \leq \dots
    \leq \lim_{n \to \infty} \B_{\pol_{k+1}}^n q_{\pol_k}
    = q_{\pol_{k+1}} .
\end{align*}

Thus we can show that MC-O-PI transitions on average to improving policies even with Tsitsiklis' framework.
Yet, \citet{Tsitsiklis2002OnIteration}, \citet{Chen2018OnProblem} and \citet{Liu2021OnStarts} do not mention this property.
% but this remark suggests that in the limit noise cannot prevent the algorithm from visiting monotonically improving policies.
\end{remark}

%%%%%%%%%%%%%%%%%%%%%%%%%%%%%%%%%%%%%%%%%
\subsection{Global convergence of mean-field solutions}

% The monotone improvement directly yields global convergence to $\qopt$.
\vspace{0.4cm}
\begin{proposition}
\label{thm: convergence deterministic mcopi}
    Solutions of the mean-field dynamics in \eqref{eq: deterministic mcopi} converge to the optimal action values $\qopt$,
    if the 
    sampling distributions $(\sigma_k)_{k \ge 0}$ satisfy exploring starts (\autoref{def: exploring starts}),
    the update distributions $(\mu_k)_{k \ge 0}$ are uniform over all actions in each state (\autoref{asp: unbiased over policies}),
    and, together with the learning rates $(\alpha_k)_{k \ge 0}$, they satisfy the modified Robbins-Monro conditions (\autoref{asp: stochastic conditions on effective learning rate}).
\end{proposition}
\vspace{-0.8cm}
\begin{proof}
Let $(\pi_k)_{k \ge 0}$ be the sequence of policies associated with a solution of \eqref{eq: deterministic mcopi}
and let $(t_n)_{n \ge 0}$ be the corresponding transition times.
By assumption, the update distributions are uniform over all actions in each state and satisfy exploring starts.
\autoref{thm: unbiased mces policy improves} thus ensures that the sequence of action values $(q_{\pi_k})_{k \ge 0}$ is componentwise non-decreasing.

Since the policy set is finite, the subsequence $(q_{\pi_{t_n}})_{n \ge 0}$ must settle after finitely many transitions.
In other words, there exists a finite index \(N \ge 0\) such that
$t_N < \infty$ and $t_{N+1} = \infty$.
Moreover, policy $\pi_{t_N}$ must be optimal. Otherwise the modified Robbins-Monro and exploring starts conditions would force $q_k$ to converge to $q_{\pi_{t_N}}$, which in turn would cause the policy to improve and thus $t_{N+1}<\infty$.


As a result, for every $k \ge t_N$, the solution of \eqref{eq: deterministic mcopi} satisfies
\[
q_{k+1} = q_k + \alpha_k \mu(\sigma_k, \best{\pi}, f_k)( \qopt - q_k) .
\]
The modified Robbins-Monro and exploring starts conditions then guarantee that $q_k$ converges to $\qopt$.
\end{proof}

% Since each $q_{\pi_k}$ belongs to the finite set $\{q_\pi : \pi \in \polset\}$, this sequence must stabilize. Thus, there exists a time step $K' \ge K$ and a policy $\pi$ such that $q_{\pi_k} = q_\pi$ for all $k \ge K'$.

% As a result, for $k \ge K'$ the solution follows the linear system
% \[
% q_{k+1} = q_k + \alpha_k \mu(\sigma_k, \pi_k, f_k)( q_\pi - q_k)
% \]
% The Robbins-Monro and exploring starts conditions in \autoref{asp: standing} guarantee that the sequence $(q_k)$ converges to $q_\pi$.

% Since $\lim_{k \to \infty} q_k = q_\pi$, \autoref{thm: asymptotic set of policies} establishes that there exists a time step $K'' \ge K'$ such that $\grpol(q_k) \subseteq \grpol(q_\pi)$ for all $k \ge K''$.
% By definition, the active policy $\pi_k$ always satisfies $\pi_k \in \grpol(q_k)$. 
% Therefore, for all $k \ge K''$
% \[
% \pi_k \in \grpol(q_\pi)
% \]
% Since $q_\pi = q_{\pi_k}$ for all $k \ge K'$, substituting this into the relation above yields $\pi_k \in \grpol(q_{\pi_k})$.
% This confirms that for all $k \ge K''$, the policy $\pi_k$ is greedy with respect to its own action values, and hence satisfies the Bellman optimality equation $q_{\pi_k} = T_{\pi_k} q_{\pi_k} = T q_{\pi_k}$.
% Consequently, $\pi_k$ is an optimal policy and the sequence $(q_k)$ converges to $q_\pi = \qopt$.

% Technically, this result does not require exploring starts. It suffices to impose Robbins-Monro learning rates and ensure that each state-action pair is updated infinitely often. We discuss this relaxation in section \ref{sec: relaxing conditions}.

% \autoref{thm: convergence deterministic mcopi} established that mean-field solutions converge to the optimal action values.
% Hence any suboptimal asymptotic behaviour of MC-O-PI must only be caused by stochastic perturbations.

The mean-field dynamics in Tsitsiklis' setting also exhibit this monotone improvement.
While previous convergence proofs cannot accommodate updates that are only uniform over actions in each state because the crucial commutation step in \eqref{eq: tsitsiklis commutation} fails,
the next section confirms that the convergence results nonetheless extend to this broader class of update distributions.


%%%%%%%%%%%%%%%%%%%%%%%%%%%%%%%%%%%%%%%%%
%%%%%%%%%%%%%%%%%%%%%%%%%%%%%%%%%%%%%%%%%
\section{MC-O-PI converges with uniform updates}

% The previous chapters showed that non-uniform update distributions can easily cause convergence to suboptimal points. For instance, if at every step the actions associated with the current greedy policy are updated significantly less frequently than non-greedy actions, stable equilibria can emerge even on the boundary of the optimal policy region.


% In contrast,
% the previous section suggests that uniform updates prevent such issues because
% the mean-field dynamics induce monotonically improving policies that drive solutions to the optimal action values $\qopt$.
% Hence any suboptimal asymptotic behaviour of MC-O-PI must only be caused by stochastic perturbations.


\autoref{thm: shrinking epsilon tube around deterministic solution} established that MC-O-PI solutions are trapped inside a shrinking tube around the reinitialised mean-field solutions,
while
Propositions \ref{thm: unbiased mces policy improves}--\ref{thm: convergence deterministic mcopi} showed that mean-field solutions converge to the optimal action values through a sequence of improving policies when updates are uniform.
This section proves that stochastic perturbations cannot consistently offset this improvement,
ensuring that MC-O-PI converges to $\qopt$ when updates are uniform over the actions of each state.


%%%%%%%%%%%%%%%%%%%%%%%%%%%%%%%%%%%%%%%%%
\subsection{Lower bound on the probability of policy improvement at transitions}

First, let us impose one additional assumption on the learning rates.

\begin{assumption}[Comparability of learning rates]
\label{asp: additional conditions on learning rate}
    The learning rates $(\lrate_k)_{k \ge 0}$ are non-increasing,
    meaning $\alpha_{k+1}\le \alpha_k$,
    and there exist a time $K_\alpha<\infty$ and a constant $\rho_\alpha\in(0,1]$ such that 
    for every $k\ge K_\alpha$
    and every $i \in \{0, \dots, \left\lfloor 1 / \alpha_k \right\rfloor\}$,
    \begin{equation}
    \label{eq:comparison}
    \alpha_{k+i}\ge \rho_\alpha \alpha_k .
    \end{equation}
\end{assumption}
\vspace{-0.6cm}

\autoref{asp: additional conditions on learning rate} is a common regularity condition when analysing asynchronous stochastic approximations
\citep{Borkar1998AsynchronousApproximations, Perkins2012AsynchronousInclusions, Liu2025TheNoise}.
It ensures the learning rates do not decay too quickly as they remain comparable to $\alpha_k$ for at least $\lfloor 1/\alpha_k \rfloor$ successive steps.
Many classical learning-rate schemes satisfy this condition, such as $\alpha_k = c_1 / (k+c_2)^{c_3}$ with $c_1>0$, $c_2 \ge 0$ and $c_3 \in (1/2, 1]$.

Consider a policy sequence $(\pi_k)_{k \ge 0}$ generated by MC-O-PI, and let $(t_n)_{n \ge 0}$ be the associated transition times.
By \autoref{def: transition times}, MC-O-PI targets the action values of policy $\pi_{t_n}$ during the time interval $\{t_n, \dots, t_{n+1}-1\}$.
We know that if $\pi_{t_n}$ is suboptimal then $t_{n+1}$ is almost surely finite, and hence $\pi_{t_{n+1}}$ is well defined.
The following result establishes a uniform lower bound on the probability that policy $\pi_{t_{n+1}}$ is better than $\pi_{t_{n}}$.

\begin{proposition}
\label{thm: fixed prob that bad actions dont become greedy}
    Let \((\pi_k)_{k\ge 0}\) be a sequence of policies generated by MC-O-PI (\autoref{def: mcopi})
    with
    learning rates $(\alpha_k)_{k \ge 0}$ that satisfy the Robbins-Monro conditions (\autoref{asp: stochastic approximation step size}) and the comparability assumption (\autoref{asp: additional conditions on learning rate}),
    sampling distributions $(\sigma_k)_{k \ge 0}$ that satisfy strict exploring starts (\autoref{def: strict exploring starts}),
    initial-visit update functions $(f_k)_{k \ge 0}$ (\autoref{def: initial visit updates}),
    and uniform update distributions $(\mu_k)_{k \ge 0}$ over all actions in each state (\autoref{asp: unbiased over policies}).
    Let \((t_n)_{n\ge 0}\) be the associated transition times (\autoref{def: transition times}).
    
    There exist a constant $p>0$ and an almost surely finite time $K$ such that
    for every $n \in \mathbb Z_{\ge 0}$ with $t_n \ge K$,
    \begin{align}
    \label{eq: porb better action value}
    \mathbb P \left(
    \pi_{t_{n+1}} \succ \pi_{t_n}
    \mid \mathcal F_{t_n}, \pi_{t_n} \notin \mathcal P_* \right) \ge p .
    \end{align}
\end{proposition}
\vspace{-0.6cm}

The main point in proving \autoref{thm: fixed prob that bad actions dont become greedy} is to construct an event with non-zero probability that ensures that
\begin{align}
\label{eq: goal of proof stoch improvement}
q_{\pi_{t_n}} ( s,\pi_{t_n} ) \le q_{\pi_{t_n}} ( s, \pi_{t_{n+1}} ) 
\end{align}
for every state $s \in \mathcal S$.
The conclusion then follows directly from \autoref{thm: strict improvement theorem}.
To construct this event, we study the evolution of $q_{k}$ in each state-action pair $(s,\pi_{t_n}(s))$ relative to the competing pairs $(s,a)$. We encode this difference using gap variables.


\begin{definition}[Gap variables]
\label{def: gap variables}
Let \((t_n)_{n \ge 0}\) be the transition times associated with the sequence of policies generated by MC-O-PI with uniform update distributions over all actions in each state (\autoref{asp: unbiased over policies}).
Fix any finite $n \in \mathbb Z_{\ge 0}$,
and let $\pi \equiv \pi_{t_n}$.
For each state $s \in \mathcal S$ and action $a \in \mathcal A(s)$, the gap is initialised at time $t_n$ as
\begin{align*}
% \label{eq: def gap initialisation}
d_{t_n}(s,a) = q_{t_n}(s,\pi) - q_{t_n}(s,a) ,
\end{align*}
and evolves for every $k \in \{t_n, \dots, t_{n+1}-1\}$ under the MC-O-PI dynamics as
\begin{align*}
% \label{eq: def gap evolution}
d_{k+1}(s,a) = d_k(s,a) + \alpha_k \big( \mu_k(s)(d_\pi(s,a) - d_k(s,a)) + \xi_k(s,a) \big)
\end{align*}
where the uniform update distributions simplify to
\begin{align*}
    \mu_k(s) \coloneqq \mu_k(s,\pi) = \mu_k(s,a) ,
\end{align*}
the target gap is
\begin{align*}
% \label{eq: def policy gap}
    d_\pi(s,a) \coloneqq q_{\pi}(s,\pi) - q_{\pi}(s,a) ,
\end{align*}
and the gap noise is
\begin{align*}
% \label{eq:xi-def}
    \xi_k(s,a) &\coloneqq
w_{k}(s,\pi) - w_{k}(s,a)
\\
\nonumber
&\: = u_k(s,\pi)v_k(s,\pi) - u_k(s,a) v_k(s,a)
\\
\nonumber
&\: \quad + \bigl( u_k(s,\pi) - \mu_k(s) \bigr) \bigl( q_\pi(s,\pi) - q_k(s,\pi) \bigr)
\\
\nonumber
&\: \quad - \bigl( u_k(s,a) - \mu_k(s) \bigr) \bigl( q_\pi(s,a) - q_k(s,a) \bigr) .
\end{align*}
\end{definition}
\vspace{-0.3cm}

% Note that we are not saying that the policy does not change between $t_n$ and $t_{n+1}-1$. Instead we just say that it does not matter because the target action values remain fixed and, as long as we can guarantee that no action from $\mathcal B_{\pi_{t_n}}$ is sampled during that time frame that is enough, i.e.
% \[
% q_{k}(s,\pi_k(s)) \ge q_{k}(s,\pi_{t_n}(s)) \ge q_{k}(s,a)
% \]


Even though $w_k(s,\pi)$ and $w_k(s,a)$ are not independent, the gap noise $\xi_k(s,a)$ is unbiased and its conditional variance is eventually bounded by a constant.


\begin{lemma}
\label{thm: properties noise gap}
For every time $k \ge 0$,
every state $s \in \mathcal S$,
and every action $a \in \mathcal A(s)$,
\begin{align}
\label{eq:xi-mean}
&\mathbb E\!\left[ \xi_{k}(s,a) \mid \mathcal F_k \right] = 0 .
\end{align}
Moreover, there exists a constant $c_\xi>0$ 
% such that for every MC-O-PI solution there is 
and an almost surely finite time $K_\xi$ such that
for every time $k \ge K_\xi$,
every state $s \in \mathcal S$,
and every action $a \in \mathcal A(s)$,
\begin{align}
\label{eq:xi-var}
&\mathbb E\!\left[ \xi_{k}^2(s,a) \mid \mathcal F_k \right] \le c_\xi .
\end{align}
\end{lemma}
\begin{proof}
\autoref{thm: properties of noise} directly implies \eqref{eq:xi-mean} since
\[
\E\!\big[ \xi_{k}(s,a) \mid \mathcal F_k \big] 
=
\E\!\big[ w_{k}(s,\pi) \mid \mathcal F_k \big] - \E\!\big[ w_{k}(s,a) \mid \mathcal F_k \big] 
= 
0 .
\]
Further,
\autoref{thm: bounded limit sets} indicates that MC-O-PI converges almost surely to the set $\bigl\{ q \in \mathcal Q : q_{\worst{\pi}} \le q \le q_{\best{\pi}}
\bigr\}$.
Thus, there exists an almost surely finite time $K$ such that $\| q_k \| \le c_{\mathcal P} + 1$ for every $k \ge K$,
with $c_{\mathcal P} := \max_{\pi \in \mathcal P}\| q_\pi \|_\infty$.
Substituting this bound into \eqref{eq: muw_second_moment} yields
\[
\mathbb E[w_k^2(s,a) \mid \mathcal F_k] 
\le c_w(1+\|q_k\|^2)
\le c_w(1+(c_{\mathcal P} + 1)^2)
\]
for every $k \ge K$, every $s \in \mathcal S$ and every $a \in \mathcal A(s)$.
Using $(x-y)^2 \le 2 ( x^2 + y^2)$ then gives
\begin{align*}   
&\mathbb E\!\left[ \bigl(w_k(s,a) - w_k(s,a') \bigr)^2 \mid \mathcal F_k \right]
\\
&\qquad\le
2 \left( \mathbb E\!\left[ w_k^2(s,a) \mid \mathcal F_k \right] 
+ \mathbb E\!\left[ w_k^2(s,a') \mid \mathcal F_k \right] \right)
\\
&\qquad\le
4 \, c_w \left( 1 + (c_{\mathcal P}+1)^2 \right) .
\end{align*}
for every $k \ge K$, every $s \in \mathcal S$, and all $a, a' \in \mathcal A(s)$.
Hence, \eqref{eq:xi-var} holds with $c_\xi := 4 \, c_w \left( 1 + (c_{\mathcal P}+1)^2 \right)$ and $K_\xi = K$.
\end{proof}


For each transition time $t_n$,
we focus on the set of actions that would worsen the policy $\pi_{t_n}$, namely
\[
\mathcal B_{\pi_{t_n}} := \big\{ (s,a) : s \in \mathcal S, \ a \in \mathcal A(s), \ q_{\pi_{t_n}}(s,a) < q_{\pi_{t_n}}(s, \pi_{t_n}) \big\} .
\]
We show that, since
(1) each gap is non-negative at initialisation because $\pi_{t_n}$ is greedy with respect to $q_{t_n}$,
and (2) the target $d_{\pi_{t_n}}$ is strictly positive for each pair in $\mathcal B_{\pi_{t_n}}$,
there is a fixed probability that no pair in $\mathcal B_{\pi_{t_n}}$ is greedy at time $t_{n+1}$, and hence that we obtain \eqref{eq: goal of proof stoch improvement}.

% At initialisation, all gaps are non-negative because $\pi_{t_n}$ is greedy with respect to $q_{t_n}$. Furthermore, the target $d_{\pi_{t_n}}$ is strictly positive for each pair in $\mathcal B_{\pi_{t_n}}$.
% We now show that, since the gap of each pair in $\mathcal B_{\pi_{t_n}}$ starts non-negative and is updated towards a strictly positive target, there is a fixed probability that no pair in $\mathcal B_{\pi_{t_n}}$ is greedy at the next transition time $t_{n+1}$, and hence obtain \eqref{eq: goal of proof stoch improvement}.


To guarantee this upwards trajectory for each pair in $\mathcal B_{\pi_{t_n}}$, we fix a margin $\delta>0$ and define three events.
The \textit{seed} event $\mathcal E^s_{t_n}$ drives each gap at least up to $O(\alpha_k)$.
The \textit{growth} event $\mathcal E^g_{t_n}$ then pushes each gap from $O(\alpha_k)$ to $\delta$.
Finally, the \textit{lock-in} event $\mathcal E^l_{t_n}$ ensures that after reaching $\delta$, each gap remains strictly positive up to time $t_{n+1}$.

Concretely, define the minimum target gap across all policies
\begin{align}
\label{eq: min target gap}  
d_{\min}
\coloneqq
\min_{\pi\in\mathcal P}\ \min_{(s,a)\in \mathcal B_\pi} d_\pi(s,a) > 0 ,
\end{align}
with the convention that $\min\varnothing=\infty$.
Using the minimum sampling probability $\sigma_{\min}$ guaranteed by the strict exploring starts,
and the comparability constant $\rho_\alpha$ from \autoref{asp: additional conditions on learning rate},
we choose $\delta$ such that
\begin{equation}
\label{eq:delta-choice}
0<\delta<\min\left\{\frac{d_{\min}}{8},\,\frac{\rho_\alpha \sigma_{\min} d_{\min}}{32}\right\}.
\end{equation}
This allows us to further define the constants
\begin{align}  
\label{eq:cs-def}
&c_s \coloneqq \frac{d_{\min}}{4} - \delta > 0,
\\
\label{eq:cg-def}
&c_g \coloneqq \sigma_{\min} ( d_{\min} - 2 \delta) > 0,
\\
\label{eq:c-def}
&c \coloneqq \min \left\{ c_s, c_g \right\} > 0.
\end{align}



%%%%%%%%%%%%%%%%%%%%%%%%%%%%%%%%%%%%%%%%%%%%%%%%%%%%%%%%%%%%%%%
\subsubsection{Seed event pushes gap up to $\bm{O(\alpha_k)}$}

Let $M \in \mathbb N$ denote a finite time horizon whose value will be fixed later.
The seed event ensures that the gap receives a guaranteed positive push of size $O(\alpha_k)$ over $M$ consecutive steps. To achieve this regardless of the current estimate $q_k$, we define the event differently depending on whether the value of the greedy action is under- or overestimated.

Concretely,
for any index $n \in \mathbb Z_{\ge 0}$,
any time $k \in \{t_n, \dots, t_{n+1}-1\}$,
any state $s \in \mathcal S$,
and any action $a \in \mathcal A(s)$,
if the value is underestimated in $s$, meaning $q_{\pi_{t_n}}(s,\pi_{t_n}) - q_{k}(s,{\pi_{t_n}}) \ge d_{\min}/2$,
we update the pair $(s,{\pi_{t_n}}(s))$ for $M$ consecutive steps and require that the noise at each step is not too negative. We gather these conditions in the event
\begin{align*}
    \mathcal E^+_k(s,a)
    \coloneqq 
    \bigcap_{k'=k}^{k+M-1}
    \left\{
    u_{k'}(s,\pi_{t_n}) = 1, \
    v_{k'}(s,\pi_{t_n}) \ge -\frac{d_{\min}}{4}
    \right\} .
\end{align*}
On the other hand, if $q_{\pi_{t_n}}(s,\pi_{t_n}) - q_{k}(s,{\pi_{t_n}}) < d_{\min}/2$,
we update $(s,a)$ and require that the noise is not too positive:
\begin{align*}
    \mathcal E^-_k(s,a)
    \coloneqq 
    \bigcap_{k'=k}^{k+M-1}
    \left\{
    u_{k'}(s,a) = 1, \
    v_{k'}(s,a) \le \frac{d_{\min}}{4}
    \right\} .
\end{align*}
Altogether, we obtain the seed event
\begin{multline}
\label{eq: def good prefix}
    \mathcal E^s_k(s,a) \coloneqq 
    \left( \left\{ q_{\pi_{t_n}}(s,\pi_{t_n}) - q_{k}(s,{\pi_{t_n}}) \ge \frac{d_{\min}}{2} \right\} \cap \mathcal E^+_k(s,a) \right) 
    \\
    \bigcup 
    \left( \left\{ q_{\pi_{t_n}}(s,\pi_{t_n}) - q_{k}(s,\pi_{t_n}) < \frac{d_{\min}}{2} \right\} \cap \mathcal E^-_k(s,a) \right) .
\end{multline}
We also define the hitting time
\[
\tau^s_k(s,a) \coloneqq \inf \big\{ k' \in \{k+1, \dots, k+M \} : q_{k'}(s,\pi_{t_n})-q_{k'}(s,a) \ge \delta \} ,
\]
and the accumulated learning rate mass between $k$ and $k'>k$ as
\[
A(k,k') \coloneqq \sum_{i=k}^{k'-1} \alpha_{i} .
\]
% with the convention that $A(k, k) = 0$.

\begin{lemma}[Seed event]
\label{thm: event to reach a_k margin}
Consider the setting of \autoref{thm: fixed prob that bad actions dont become greedy}.
For every index $n \in \mathbb Z_{\ge 0}$,
every pair $(s,a) \in \mathcal B_{\pi_{t_n}}$
and every time $k \in \{t_n, \dots, t_{n+1}-1\}$,
the joint event
\begin{align}
\label{eq: joint seed event}
\mathcal E^s_k(s,a)
\cap
\big\{0 \leq q_{k}(s,\pi_{t_n}) - q_{k}(s,a) < \delta \big\}
\end{align}
% guarantees that 
ensures
% for every $k'$ that satisfies $k < k' \le \big( \tau^s_k(s,a) \wedge (k+M) \wedge t_{n+1} \big)$,
for every $k'$ that satisfies $k < k' \le \min\{ \tau^s_k(s,a), k+M, t_{n+1} \}$,
% for every $k' \in \{k +1, \dots, \min\{ \tau^s_k(s,a), k+M, t_{n+1} \}\}$,
\begin{align}
\label{eq: lower bound during good prefix}
    q_{k'}(s,\pi_{t_n}) - q_{k'}(s,a)
    > c_s A(k,k') > 0 .
\end{align}
% for every $k'$ that satisfies $k < k' \le \big( \tau^s_k(s,a) \wedge (k+M) \wedge t_{n+1} \big)$.
Moreover, the probability of $\mathcal E^s_k(s,a)$ is lower bounded by
\begin{align}
\label{eq: probability reach a_k margin}
    \mathbb P ( \mathcal E^s_k(s,a) \mid \mathcal F_k, \pi_{t_n} ) \ge \left( \sigma_{\min} \frac{d_{\min}^2}{d_{\min}^2 + 16 c_v} \right)^M > 0 .
\end{align}
\end{lemma}

\begin{proof}
Fix an $n \in \mathbb Z_{\ge 0}$ and let $\pi \equiv \pi_{t_n}$.
Fix a pair $(s,a) \in \mathcal B_\pi$,
define the sequence $\big( d_k(s,a) \big)_{k=t_n}^{t_{n+1}}$ as in \autoref{def: gap variables}, and write $d_k \equiv d_k(s,a)$ and $d_\pi \equiv d_\pi(s,a)$ for brevity.
Fix a time $k \in \{t_n, \dots, t_{n+1}-1\}$ and write $\tau^s_k \equiv \tau^s_k(s,a)$.
We split the proof along the two mutually exclusive cases that make up $\mathcal E^s_k(s,a)$.


First consider the case when $q_\pi(s,\pi) - q_{k}(s,\pi) \ge d_{\min}/2$.
On $\mathcal E^+_k(s,a)$, only $(s, \pi(s))$ is updated during $M$ consecutive steps,
and $v_{k'}(s,\pi) \ge -d_{\min}/4$ for every $k' \in \{k, \dots, k+M-1\}$.
Hence, for every $k'$ that satisfies 
% $k \le k' < \big((k+M) \wedge t_{n+1} \big)$,
$k \le k' < \min\{k+M, t_{n+1} \}$,
\begin{align}
\nonumber
    d_{k'+1}
    &= d_{k'} + \alpha_{k'} \big( q_\pi(s,\pi) - q_{k'}(s,\pi) + v_{k'}(s,\pi) \big)
    \\
\label{eq: plug noise into update of seed stage case 1}
    &\ge d_{k'} + \alpha_{k'} \left( q_\pi(s,\pi) - q_{k'}(s,\pi) -\frac{d_{\min}}{4} \right) .
\end{align}
Since $d_{k} \ge 0$ on $\{0 \le d_{k} < \delta \}$
and $d_{k'} < \delta$ before the hitting time $\tau^s_k$,
it follows that for every $k'$ that satisfies 
% $k \le k' < \big( \tau^s_k \wedge (k+M) \wedge t_{n+1} \big)$,
$k \le k' < \min \{ \tau^s_k, k+M, t_{n+1} \}$,
\begin{align}
\nonumber
    q_\pi(s,\pi) - q_{k'}(s,\pi)
    &= q_\pi(s,\pi) - d_{k'} + d_{k} - q_{k}(s,\pi)
    \\
\label{eq: lower bound update step in seed stage case 1}
    &> \frac{d_{\min}}{2} - \delta .
\end{align}
Equations \eqref{eq: plug noise into update of seed stage case 1} and \eqref{eq: lower bound update step in seed stage case 1}
indicate that on the joint event \eqref{eq: joint seed event}
the gap strictly increases at each time $k'$ in the restricted range:
\begin{align}
\label{eq: final lower bound case 1}
    d_{k'+1}
    > d_{k'} + \alpha_{k'} \left( \frac{d_{\min}}{4} - \delta \right)
    = d_{k'} + \alpha_{k'} c_s .
    % > 0 .
\end{align}

Now consider the case when $q_\pi(s,\pi) - q_{k}(s,\pi) < d_{\min}/2$.
On $\mathcal E^-_k(s,a)$, only $(s, a)$ is updated during $M$ consecutive steps, and $v_{k'}(s,a) \le d_{\min}/4$ for every $k' \in \{k, \dots, k+M-1\}$.
Hence, for every $k'$ that satisfies 
% $k \le k' < \big((k+M) \wedge t_{n+1} \big)$,
$k \le k' < \min \{k+M, t_{n+1} \}$,
\begin{align}
\nonumber
    d_{k'+1}
    &= d_{k'} + \alpha_{k'}(q_{k'}(s,a) - q_\pi(s,a) - v_{k'}(s,a))
    \\
\label{eq: plug noise into update of seed stage case 2}
    &\ge d_{k'} + \alpha_{k'} \left( q_{k'}(s,a) - q_\pi(s,a) - \frac{d_{\min}}{4} \right) .
\end{align}
Since $q_{k}(s,\pi)$ is unchanged on $\mathcal E^-_k(s,a)$
and $d_{k'} < \delta$ before $\tau^s_k$,
it follows that, for every $k'$ that satisfies 
% $k \le k' < \big( \tau^s_k \wedge (k+M) \wedge t_{n+1} \big)$,
$k \le k' < \min\{ \tau^s_k , k+M , t_{n+1} \}$,
\begin{align}
\nonumber
    q_{k'}(s,a) - q_\pi(s,a)
    &= d_\pi - ( q_\pi(s,\pi) - q_{k}(s,\pi) ) - d_{k'}
    \\
\label{eq: lower bound update step in seed stage case 2}
    &> d_{\min} - \frac{d_{\min}}{2} - \delta
    \ge \frac{d_{\min}}{2} - \delta .
\end{align}
Equations \eqref{eq: plug noise into update of seed stage case 2} and \eqref{eq: lower bound update step in seed stage case 2}
confirm that on the joint event \eqref{eq: joint seed event}
the gap strictly increases at each time $k'$ in the restricted range:
\begin{align}
\label{eq: final lower bound case 2}
    d_{k'+1}
    > d_{k'} + \alpha_{k'} \left( \frac{d_{\min}}{4} - \delta \right)
    = d_{k'} + \alpha_{k'} c_s .
    % > 0 .
\end{align}

By induction on \eqref{eq: final lower bound case 1} or \eqref{eq: final lower bound case 2} and shifting the index by 1 to include time $k'+1$, we obtain
\begin{align}
\label{eq: increase in d^b case B}
    d_{k'}
    > d_{k} + c_s A(k,k')
    \ge c_s A(k,k') > 0
\end{align}
as long as 
% $k < k' \le \big( \tau^s_k \wedge (k+M) \wedge t_{n+1} \big)$, 
$k < k' \le \min\{ \tau^s_k , k+M , t_{n+1} \}$, 
which proves \eqref{eq: lower bound during good prefix}.

Finally, we lower-bound the probability of $\mathcal E^s_k(s,a)$.
Given that $v_k$ is unbiased and its variance is bounded by $c_v$ uniformly over the policies,
Cantelli's inequality implies that for every $k' \ge 0$, $s \in \mathcal S$ and $a \in \mathcal A(s)$,
\begin{align}
&\mathbb P \! \left( v_{k'}(s,a) \ge -\frac{d_{\min}}{4} \,\middle|\, \mathcal F_{k'}, u_{k'}(s,a)=1 \right) \ge \frac{d_{\min}^2}{d_{\min}^2 + 16 c_v},
\label{eq:noise-lower}
\\
&\mathbb P \! \left( v_{k'}(s,a) \le \frac{d_{\min}}{4} \,\middle|\, \mathcal F_{k'}, u_{k'}(s,a)=1 \right) \ge \frac{d_{\min}^2}{d_{\min}^2 + 16 c_v}.
\label{eq:noise-upper}
\end{align}
Since the case choice in \eqref{eq: def good prefix} is measurable on $(\mathcal F_{k}, \pi_{t_n})$ and the minimum update probability of any state-action pair is $\sigma_{\min}$ under strict exploring starts, 
repeated conditioning over $M$ steps yields
\begin{align}
\label{eq: prob of good prefix}
\mathbb P \big( \mathcal E^s_k(s,a) \mid \mathcal F_{k} , \pi_{t_n} \big) \ge \left( \sigma_{\min} \frac{d_{\min}^2}{d_{\min}^2 + 16 c_v} \right)^M > 0 .
\end{align}

Since $n$, $(s,a)$ and $k$ were chosen arbitrarily from their respective sets, the result holds uniformly.
\end{proof}

%%%%%%%%%%%%%%%%%%%%%%%%%%%%%%%%%%%%%%%%%%%%%%%%%%%%%%%%%%%%%%%
\subsubsection{Growth event pushes gap from $\bm{O(\alpha_k)}$ to $\bm{\delta}$}

% To measure how long a gap takes to grow up to $\delta$, we 
Let $N(k)$ denote the minimum number of steps required to accumulate a particular amount of learning rate mass:
\begin{align}
N(k) 
\coloneqq
\inf \left\{ i \ge 1 : 
A(k,k+i)
\ge \frac{4 \delta}{c_g} \right\} .
\end{align}

\begin{lemma}
For every sequence of learning rates $(\alpha_k)_{k \ge 0}$ that satisfies the Robbins-Monro conditions (\autoref{asp: stochastic approximation step size}) and the comparability assumption (\autoref{asp: additional conditions on learning rate}),
there exists a time $K_N < \infty$ such that for every $k \ge K_N$,
\[N(k) \le \frac{1}{\alpha_{k}} .\]
\end{lemma}
\vspace{-0.5cm}
\begin{proof}
The definitions of $\delta$ in \eqref{eq:delta-choice} and $c_g$ in \eqref{eq:cg-def} ensure that
\begin{multline*}
    \frac{4 \delta}{c_g \rho_\alpha}
    < 
    \frac{4 \rho_\alpha \sigma_{\min} d_{\min}}{32 \sigma_{\min}(d_{\min} - 2 \delta) \rho_\alpha}
    =
    \frac{d_{\min}}{8 (d_{\min} - 2 \delta)}
    \\
    <
    \frac{d_{\min}}{8(d_{\min} - d_{\min}/4)}
    =
    \frac{1}{6}
    < 1 .
\end{multline*}
Since $(\alpha_k)_{k \ge 0}$ decreases to 0 by the Robbins-Monro conditions, there exists a finite $K$ such that $4\delta / c_g \rho_\alpha + \alpha_k < 1$ for all $k \ge K$,
and hence
\[
\frac{4 \delta}{c_g \rho_\alpha \alpha_k} + 1 < \frac{1}{\alpha_k} .
\]
Let
\[
r_k := \left\lceil \frac{4 \delta}{c_g \rho_\alpha \alpha_k} \right\rceil .
\]
Since $r_k \le 1/ \alpha_k$,
the comparability property of the learning rates guarantees that $A(k, k+r_{k}) \ge r_{k} \rho_\alpha \alpha_{k} \ge 4 \delta / c_g$ for every $k \ge \max\{K, K_\alpha\}$.
% \[
% A(k, k+r_{k}) \ge r_{k} \rho_\alpha \alpha_{k} \ge 4 \delta / c_g .
% \]
As a result, there exists a finite time $K_N \ge \max\{K, K_\alpha\}$ such that $N(k) \le r_k \le 1/\alpha_{k}$ for all $k \ge K_N$.
\end{proof}

To ensure that stochastic perturbations do not overwhelm the mean-field drift, we bound the perturbations.
For any index $n \in \mathbb Z_{\ge 0}$,
any time $k \in \{t_n, \dots, t_{n+1}-1\}$,
any state $s \in \mathcal S$
and any action $a \in \mathcal A(s)$,
define the accumulated perturbations between $k$ and $k'>k$ as
\[
\Xi(k, k' \mid s,a) 
\coloneqq 
% \sum_{i=k}^{k'-1} \alpha_{i} \xi_i(s,a)
% =
\sum_{i=k}^{k'-1} \alpha_{i} \big( w_i(s, \pi_{t_n})-w_i(s,a) \big) 
\]
with the convention $\Xi(k, k \mid s,a) = 0$.
We then define the growth event
\[
\mathcal E^g_k(s,a)
\coloneqq
\bigcap_{k'=k+1}^{k+N(k)}
\left\{
\Xi(k,k' \mid s,a) > - \frac{c_s}{2} \rho_\alpha M \alpha_k - \frac{c_g}{2} A(k,k')
\right\}
\]
for any $M \in \mathbb N$,
and define the hitting time
\[
\tau^g_k(s,a) \coloneqq \inf \! \big\{ k' \in \{k+1, \dots, k+N(k) \} : q_{k'}(s,\pi_{t_n}) - q_{k'}(s,a) \ge \delta \big\} .
\]

\vspace{-0.4cm}

\begin{lemma}[Growth event]
\label{thm: from a_k margin to delta}
Consider the setting of \autoref{thm: fixed prob that bad actions dont become greedy}.
For every index $n \in \mathbb Z_{\ge 0}$ such that $t_n \ge K_N$,
every pair $(s,a) \in \mathcal B_{\pi_{t_n}}$
and every time $k \in \{t_n, \dots, t_{n+1}-1\}$,
the joint event
\begin{align}
\label{eq: joint growth event}
\mathcal E^g_k(s,a) \cap \big\{ c_s \rho_\alpha M \alpha_k \leq q_{k}(s,\pi_{t_n}) - q_{k}(s,a) < \delta \big\}
\end{align}
guarantees that the gap hits $\delta$ within $N(k)$ steps:
\begin{align}
\label{eq: hit time event a_k to delta}
    % \big( \tau^g_k(s,a) \wedge t_{n+1} \big) \le k + N(k) ,
    \min\{ \tau^g_k(s,a) , t_{n+1} \} \le k + N(k) ,
\end{align}
while for every $k'$ that satisfies 
% $k \le k' \le \big( \tau^g_k(s,a) \wedge t_{n+1} \big)$,
$k \le k' \le \min\{ \tau^g_k(s,a), t_{n+1} \}$,
\begin{align}
\label{eq: growth event gap remains positive}
    q_{k'}(s,\pi_{t_n}) - q_{k'}(s,a) > 0 .
\end{align}
Moreover, there exists an almost surely finite time $K_\xi$ such that for every $k \ge K_\xi$,
\begin{align}
\label{eq: prob from a_k reach delta}
    \mathbb P (\mathcal E^g_k(s,a) \mid \mathcal F_k, \pi_{t_n})
    \ge 1 - \frac{8 c_\xi}{c^2 \rho_\alpha^2} \sum_{j = 0}^{\infty} \frac{2^j}{(M + 2^j)^2} .
\end{align}
\end{lemma}
\begin{proof}
Fix an $n \in \mathbb Z_{\ge 0}$ such that $t_n \ge K_N$,
and let $\pi \equiv \pi_{t_n}$.
Fix a pair $(s,a) \in \mathcal B_\pi$,
define the sequence $\big( d_k(s,a) \big)_{k=t_n}^{t_{n+1}}$ as in \autoref{def: gap variables}, and write $d_k \equiv d_k(s,a)$, $d_\pi \equiv d_\pi(s,a)$ and $\xi_k \equiv \xi_k(s,a)$ for brevity.
Fix a time $k \in \{t_n, \dots, t_{n+1}-1\}$ and write $\tau^g_k \equiv \tau^g_k(s,a)$.

For every $k' \in \{k, \dots, t_{n+1}-1\}$,
the gap satisfies
\[
d_{k'+1}
=
d_{k'} + \alpha_{k'} \big( \mu_{k'}(s) (d_\pi - d_{k'}) + \xi_{k'} \big) ,
\]
where the definition of $\delta$ guarantees that $d_\pi \ge d_{\min} > 2 \delta$,
and the strict exploring starts ensure that $\mu_{k'}(s) \ge \sigma_{\min}$. 
Since $d_{k'} < \delta$ before the hitting time $\tau^g_k$,
we can bound the update for all 
% $k \le k' < ( \tau^g_k \wedge t_{n+1} )$ 
$k \le k' < \min\{ \tau^g_k, t_{n+1} \}$ 
as
\begin{align*}
d_{k'+1}
>
d_{k'} + \alpha_{k'} \left( \sigma_{\min}(d_{\min} - \delta) + \xi_{k'} \right)
% \\
% &>
% d_{k'} + \alpha_{k'} \left( \sigma_{\min}(d_{\min} - 2 \delta) + \xi_{k'} \right)
% \\
>
d_{k'} + \alpha_{k'} \left( c_g + \xi_{k'} \right) ,
\end{align*}
using $c_g \coloneqq \sigma_{\min} ( d_{\min} - 2 \delta)$.
Summing from $k$ to $k'$ then yields
\[
d_{k'}
>
d_{k} + c_g A(k,k') + \Xi(k,k' \mid s,a) .
\]
Therefore, on the joint event $\mathcal E^g_k(s,a) \cap \{c_s \rho_\alpha M \alpha_k \leq d_k < \delta \}$,
% the gap satisfies
\begin{align}
d_{k'}
&>
c_s \rho_\alpha M \alpha_k + c_g A(k,k') - \frac{c_s}{2} \rho_\alpha M \alpha_k - \frac{c_g}{2} A(k,k')
\nonumber
\\
\label{eq: gap remains positive up to tau'_k}
&=
\frac{c_s}{2} \rho_\alpha M \alpha_k + \frac{c_g}{2} A(k,k') > 0
\end{align}
as long as 
% $k \le k' \le \big( \tau^g_k \wedge t_{n+1} \wedge (k+N(k))\big)$, 
$k \le k' \le \min \{ \tau^g_k , t_{n+1} , k+N(k) \}$.


Now assume that $t_{n+1} > k+N(k)$, otherwise 
% $(\tau^g_k \wedge t_{n+1}) \le k+N(k)$ 
$\min\{ \tau^g_k , t_{n+1} \} \le k+N(k)$ 
is trivial.
Suppose for contradiction that $\tau^g_k = \infty$ because $d_k$ does not reach $\delta$ within $N(k)$ steps.
Equation \eqref{eq: gap remains positive up to tau'_k} implies that
\[
d_{k+N(k)}
>
\frac{c_s}{2} \rho_\alpha M \alpha_k + \frac{c_g}{2} A(k,k+N(k)) .
\]
By definition, however, $A(k,k+N(k)) \ge 4 \delta/c_g$. So we obtain a contradiction: $d_{k+N(k)}
>
c_s \rho_\alpha M \alpha_k / 2 + 2 \delta
> \delta .$
% \[
% d_{k+N(k)}
% >
% \frac{c_s}{2} \rho_\alpha M \alpha_k + 2 \delta
% > \delta .
% \]
Thus, $\tau^g_k \le k+N(k)$, which proves \eqref{eq: hit time event a_k to delta}
and, together with \eqref{eq: gap remains positive up to tau'_k}, proves \eqref{eq: growth event gap remains positive}.

Finally, we lower-bound the probability of $\mathcal E^g_k(s,a)$.
We observe that the complement event $\overline{\mathcal E^g_k(s,a)}$ implies that for some $i \in \{1, \dots , N(k) \}$,
\[
\bigl| \Xi(k,k+i \mid s,a) \bigr|
\ge \frac{c_s}{2} \rho_\alpha M \alpha_k + \frac{c_g}{2} A(k,k+i) .
\]
Using $c \coloneqq \min\{c_s,c_g\}$ and the comparability property (since $i \le N(k) \le 1/\alpha_k$ for $k \ge t_n \ge K_N$), this simplifies to
\[
\bigl| \Xi(k,k+i \mid s,a) \bigr|
\ge \frac{c}{2} \big( \rho_\alpha M \alpha_k + A(k,k+i) \big)
\ge \frac{c}{2} \rho_\alpha \alpha_k ( M + i ).
\]
Choosing $j \in \mathbb N$ such that $2^j \le i \le 2^{j+1}$ yields the bound
\begin{align}
    \max_{1 \le \ell \le \min\{ 2^{j+1} , N(k)\}}
    |\Xi(k,k + \ell \mid s,a)|
    \ge \frac{c}{2} \rho_\alpha \alpha_k (M + 2^j) .
\end{align}
By \autoref{thm: properties noise gap}, there exists an almost surely finite $K_\xi$ such that $\big(\Xi(k, k+\ell \mid s,a)\big)_{\ell \ge 1}$ is a square-integrable martingale adapted to the filtration $(\mathcal F_{k+\ell})_{\ell \ge 1}$ for every $k \ge K_\xi$. Therefore, by Doob's martingale inequality,
\begin{align}
\label{eq: prob of not E^g_k}
    &\mathbb P
    \left( 
    \overline{\mathcal E^g_k(s,a)} \mid \mathcal F_{k}, \pi_{t_n}
    \right)
    \\
    \nonumber
    &\qquad \le
    \sum_{j = 0}^{\infty} \frac{\mathbb E \left[ \big( \Xi(k, \min\{ k+2^{j+1} , k+N(k) \} \mid s,a) \big)^2 \mid \mathcal F_{k}, \pi_{t_n} \right]}{\left( \dfrac{c}{2} \rho_\alpha \alpha_k ( M + 2^j) \right)^2} .
\end{align}
\autoref{thm: properties noise gap} and the non-increasing sequence $(\alpha_k)_{k \ge 0}$ imply that
\begin{multline*}
\mathbb E \left[ \big( \Xi(k, \min \{ k+2^{j+1} , k+N(k) \} \mid s,a) \big)^2 \mid \mathcal F_{k}, \pi_{t_n} \right]
\\
\le
c_\xi \sum_{\ell=0}^{2^{j+1}-1} \alpha^2_{k+\ell}
\le
c_\xi 2^{j+1} \alpha^2_k
\end{multline*}
for every $k \ge K_\xi$.
Substituting this back into \eqref{eq: prob of not E^g_k} proves \eqref{eq: prob from a_k reach delta}:
\begin{align}
    \mathbb P
    \left( 
    \overline{\mathcal E^g_k(s,a)} \mid \mathcal F_{k}, \pi_{t_n}
    \right)
    \le
    \frac{8 c_\xi}{c^2 \rho_\alpha^2} \sum_{j = 0}^{\infty} \frac{2^j}{(M+ 2^j)^2}.
\end{align}

Since $n$, $(s,a)$ and $k$ were chosen arbitrarily from their respective sets, the result holds uniformly.
\end{proof}

% \autoref{thm: from a_k margin to delta} confirms that on the event $\mathcal E^g_k \cap \{c \rho_\alpha M \alpha_k \leq q_k(s,\pi_{t_n}(s))-q_k(s,a) < \delta \}$, the gap safely hits $\delta$ in at most $N(k)$ steps while remaining strictly positive throughout the ascent.

%%%%%%%%%%%%%%%%%%%%%%%%%%%%%%%%%%%%%%%%%%%%%%%%%%%%%%%%
\subsubsection{Lock-in event keeps gap positive after reaching $\bm{\delta}$}

Once the gap reaches the safe margin $\delta$, the lock-in event ensures that the remaining weighted stochastic perturbations are too small to drag the gap below zero.
For any $n \in \mathbb Z_{\ge 0}$,
any state $s \in \mathcal S$
and any action $a \in \mathcal A(s)$,
we define the lock-in event as
\begin{align}
    \mathcal E^l_{t_n}(s,a) \coloneqq \left\{ \sup_{k \ge t_n} \bigl| \Xi(t_n, k \mid s,a) \bigr| < \frac{\delta}{8} \right\} .
\end{align}

\begin{lemma}[Lock-in event]
\label{thm: event noise remains small}
Consider the setting of \autoref{thm: fixed prob that bad actions dont become greedy}.
For every index $n \in \mathbb Z_{\ge 0}$,
every pair $(s,a) \in \mathcal B_{\pi_{t_n}}$
and every time $k \in \{t_n, \dots, t_{n+1}-1\}$,
the joint event
\[
\mathcal E^l_{t_n}(s,a) \cap \big\{ q_{k}(s,\pi_{t_n}) - q_{k}(s,a) \ge \delta \big\}
\]
guarantees that
% the gap remains strictly positive until the next transition time, meaning
for every $k' \in \{k, \dots, t_{n+1}\}$,
\begin{align}
\label{eq: gap remains strictly positive}
q_{k'}(s,\pi_{t_n}) > q_{k'}(s,a) .
\end{align}
Moreover, there exists an almost surely finite time $K_\xi$ such that for every $t_n \ge K_\xi$,
\begin{equation}
\label{eq:probability-small-remaining-noise}
\mathbb P\!\left(\mathcal E^l_{t_n}(s,a)\mid \mathcal F_{t_n}, \pi_{t_n} \right)
\ge
1-\frac{64 c_\xi}{\delta^2}\sum_{i=t_n}^\infty \alpha_i^2 .
\end{equation}
\end{lemma}
\vspace{-0.6cm}
\begin{proof}
Fix an $n \in \mathbb Z_{\ge 0}$
and let $\pi \equiv \pi_{t_n}$.
Fix a pair $(s,a) \in \mathcal B_\pi$,
define the sequence $\big( d_k(s,a) \big)_{k=t_n}^{t_{n+1}}$ as in \autoref{def: gap variables}, and write $d_k \equiv d_k(s,a)$, $d_\pi \equiv d_\pi(s,a)$ and $\xi_k \equiv \xi_k(s,a)$ for brevity.

For every time $k \in \{t_n, \dots, t_{n+1}-1\}$, the gap is updated as
\begin{align*}
d_{k+1}
=
\big( 1 - \alpha_{k} \mu_{k}(s) \big) d_{k} + \alpha_{k} \mu_{k}(s) d_\pi + \alpha_{k} \xi_{k} .
\end{align*}
To unroll this recursion, define the contraction factor between steps $k$ and $k'>k$ as
\[
\Gamma(k,k') := \prod_{i=k}^{k'-1} \big( 1 - \alpha_{i} \mu_{i}(s) \big) \in(0,1]
\]
with the convention that $\Gamma(k,k) = 1$.
Fix a time $k \in \{t_n, \dots, t_{n+1}-1\}$.
By induction, we obtain for every $k' \in \{k,\dots,t_{n+1}\}$,
\begin{align}
\label{eq: induction for abel proof}  
d_{k'}
=
\Gamma(k,k') d_{k}
+
\bigl(1-\Gamma(k,k') \bigr) d_\pi
+ \sum_{i=k}^{k'-1} \Gamma(i+1,k') \alpha_i \xi_i .
\end{align}
Since $d_\pi \ge d_{\min} > \delta$ and $\Gamma(k,k') \in (0,1]$, the event $\{d_{k} \ge \delta\}$ ensures that
\begin{align}
\label{eq: lower obund for abel proof}
\Gamma(k,k') d_{k} + \big(1-\Gamma(k,k')\big) d_\pi > \delta.
\end{align}
To bound the stochastic part, we note that for a fixed terminal index $k'$, the map $i \mapsto \Gamma(i+1,k')$ is nondecreasing with values in $(0,1]$.
Hence applying summation by parts and the triangle inequality yields
\begin{align}
\nonumber
&\left|
\sum_{i=k}^{k'-1} \Gamma(i+1,k') \alpha_i \xi_i
\right|
\\
\nonumber
&\qquad\le
\big| \Gamma(k',k') \big| \big| \Xi(k,k') \big| + \sum_{i=k}^{k'-2} \big| \Gamma(i+2,k') - \Gamma(i+1,k') \big| \big| \Xi(k, i+1) \big|
\\
\nonumber
&\qquad\le
\left( \Gamma(k',k') + \sum_{i=k}^{k'-2} \big( \Gamma(i+2,k') - \Gamma(i+1,k') \big) \right)
\sup_{k \le m \le k'} \bigl| \Xi(k,m\mid s,a) \bigr|
\\
\nonumber
&\qquad=
\big( \Gamma(k',k') + \Gamma(k',k') - \Gamma(k+1,k') \big)
\sup_{k \le m \le k'} \bigl| \Xi(k,m\mid s,a) \bigr|
\\
% \label{eq:abel equation on noise}
\nonumber
&\qquad\le
2 \sup_{k \le m \le k'} \bigl| \Xi(k,m\mid s,a) \bigr| .
\end{align}
On $\mathcal E^l_{t_n}(s,a)$, the triangle inequality implies that for every $m > k$,
\begin{align*}
\bigl|\Xi(k,m\mid s,a)\bigr|
&=
\bigl|
\Xi(t_n , m \mid s,a) - \Xi(t_n, k \mid s,a)
\bigr|
\\
&\le
\bigl|\Xi(t_n,m\mid s,a)\bigr|
+
\bigl|\Xi(t_n,k\mid s,a)\bigr|
\\
&<
\frac{\delta}{4}.
\end{align*}
This strictly limits the noise term in \eqref{eq: induction for abel proof} by
\begin{equation}
\label{eq:shifted-noise-bound}
\left|
\sum_{i=k}^{k'-1} \Gamma(i+1,k') \alpha_i \xi_{i}
\right|
< \frac{\delta}{2}.
\end{equation}
Plugging \eqref{eq: lower obund for abel proof} and \eqref{eq:shifted-noise-bound} into \eqref{eq: induction for abel proof}  
yields
\[
d_{k'}(s,a)
>
\delta - \frac{\delta}{2}
= \frac{\delta}{2}
>
0
\]
for every $k' \in \{k,\dots,t_{n+1}\}$, which proves \eqref{eq: gap remains strictly positive}.


Finally, we lower-bound the probability of $\mathcal E^l_{t_n}(s,a)$.
We observe that the complement event $\overline{\mathcal E^l_{t_n}(s,a)}$ implies that for some time $k \ge t_{n}$,
\[
\bigl| \Xi(t_n, k \mid s,a)\bigr|
\ge \frac{\delta}{8} ,
\]
or equivalently
\[
% \left\{\sup_{k\ge t_n}|\Xi(t_n,k\mid s,a)|\ge \frac{\delta}{8}\right\}
\overline{\mathcal E^l_{t_n}(s,a)}
=
\bigcup_{t = t_n}^\infty
\left\{
\max_{t_n\le k\le t}|\Xi(t_n,k\mid s,a)|
\ge \frac{\delta}{8}
\right\} .
\]
Thus the monotone convergence of conditional probabilities implies that
\begin{align}
% \nonumber
\label{eq:hold-doob-trunc}
\mathbb P\!\left(
\overline{\mathcal E^l_{t_n}(s,a)}
\,\middle|\,
\mathcal F_{t_n},\pi_{t_n}
\right)
&=
\lim_{t \to \infty}
\mathbb P\!\left(
\max_{t_n \le k \le t}| \Xi(t_n, k\mid s,a) |
\ge \frac{\delta}{8}
\,\middle|\,
\mathcal F_{t_n},\pi_{t_n}
\right) .
% \\
% \label{eq:hold-doob-trunc}
% &\le
% \frac{64}{\delta^2}
% \limsup_{m\to\infty}
% \mathbb E\!\left[M_m^2\mid \mathcal F_{t_n},\pi_{t_n}\right].
\end{align}
\autoref{thm: properties noise gap} indicates that there exists an almost surely finite time $K_\xi$ such that
$\big(\Xi(t_n, k \mid s,a)\big)_{t_n\le k \le t}$ is a square-integrable martingale adapted to the filtration $(\mathcal F_{k})_{t_n \le k\le t}$ for all $t \ge t_n \ge K_\xi$.
Hence, Doob's $L^2$ inequality gives
\begin{align}
\label{eq: doob intermediary}    
\mathbb P\!\left(
\max_{t_n \le k \le t} \bigl| \Xi(t_n, k \mid s,a)\bigr|
\ge \frac{\delta}{8}
\,\middle|\,
\mathcal F_{t_n}, \pi_{t_n}
\right)
\le
\frac{64}{\delta^2}
\mathbb E\!\left[\Xi(t_n,t\mid s,a)^2\mid \mathcal F_{t_n}, \pi_{t_n} \right].
\end{align}
% Using the orthogonality of martingale increments and the variance bound from \eqref{eq:xi-var}, we obtain
By using the tower property and \autoref{thm: properties noise gap}, we then obtain
\begin{align*}
\mathbb E\!\left[\Xi(t_n, t \mid s,a)^2\mid \mathcal F_{t_n}, \pi_{t_n} \right]
&=
\sum_{i=t_n}^{t-1} \alpha_{i}^2\,
\mathbb E\!\left[\xi_{i}^2\mid \mathcal F_{t_n}, \pi_{t_n} \right]
\\
&=
\sum_{i=t_n}^{t-1} \alpha_{i}^2\,
\mathbb E\!\left[
\mathbb E\!\left[\xi_{i}^2\mid \mathcal F_{i}, \pi_{t_n}\right]
\middle|\mathcal F_{t_n} , \pi_{t_n}
\right]
\\
&\le
c_\xi \sum_{i=t_n}^{t-1} \alpha_{i}^2
% \\
% &\le
% c_\xi \sum_{i=t_n}^{\infty} \alpha_{i}^2 .
\end{align*} 
Substituting this back into \eqref{eq: doob intermediary} then \eqref{eq:hold-doob-trunc} proves \eqref{eq:probability-small-remaining-noise} because
\[
\mathbb P\!\left(
\overline{\mathcal E^l_{t_n}(s,a)}
\,\middle|\,
\mathcal F_{t_n}, \pi_{t_n}
\right)
\le
\frac{64 c_\xi}{\delta^2}\sum_{i=t_n}^{\infty} \alpha_{i}^2 .
\]

Since $n$, $(s,a)$, and $k$ were chosen arbitrarily from their respective sets, the result holds uniformly.
\end{proof}

% In other words, for sufficiently large $k$, whenever a gap $d^b_k$ reaches a fixed $\delta > 0$, it remains positive with large probability. 


%%%%%%%%%%%%%%%%%%%%%%%%%%%%%%%%%%%%%%%%%%%%%%%%%%%%%%%%%%%%%%%
\subsubsection{Event for all gaps to remain strictly positive}

We now apply the three sequential events to every pair in $\mathcal B_{\pi_{t_n}}$ so that none of them become greedy up to and including the transition time $t_{n+1}$.
Fix an $n \in \mathbb Z_{\ge 0}$,
let $L := |\mathcal B_{\pi_{t_n}}|$,
and assume that $L>0$, otherwise \autoref{thm: fixed prob that bad actions dont become greedy} is trivial because any transition would improve the policy (\autoref{thm: strict improvement theorem}).
We first process all pairs in $\mathcal B_{\pi_{t_n}}$ by sequentially applying the $M$-step seed block from \autoref{thm: event to reach a_k margin}.
To do so, we define staggered start times
\[
b_\ell := t_n + (\ell-1) M
% \qquad \forall \, ,
\]
for $\ell \in \{1, \dots, L\}$,
along with an overall finish time for the seed phase
\[
T_n := t_n + L M.
\]
At time $T_n$, we then apply the $N(T_n)$-step growth block 
of \autoref{thm: from a_k margin to delta} 
to guarantee that all gaps reach the $\delta$-margin,
followed by the global lock-in event
of \autoref{thm: event noise remains small} 
to keep them strictly positive until time $t_{n+1}$.


Conditioned on $(\mathcal F_{t_n}, \pi_{t_n})$, we can enumerate the bad state-action pairs as
\[
\mathcal B_{\pi_{t_n}} \equiv \big\{ (s_1,a_1), \dots, (s_L,a_L) \big\} ,
\]
and then define the global seed, growth and lock-in events as the intersections of their individual events:
\begin{align}
\label{eq: def global seed event}
\mathcal E^s_{t_n} &\coloneqq \bigcap_{\ell=1}^{L} \mathcal E^s_{b_\ell}(s_\ell, a_\ell) ,
\\
\label{eq: def global grow event}
\mathcal E^g_{t_n} &\coloneqq \bigcap_{\ell=1}^{L} \mathcal E^g_{T_n}(s_\ell, a_\ell) ,
\\
\label{eq: def global hold event}
\mathcal E^l_{t_n} &\coloneqq \bigcap_{\ell=1}^{L} \mathcal E^l_{t_n}(s_\ell, a_\ell) .
\end{align}


Since we process the pairs in $\mathcal B_{\pi_{t_n}}$ sequentially, the gap for a specific pair $(s_\ell, a_\ell)$ might remain zero while it waits for its dedicated seed block to begin at time $b_\ell$. 
However, the policy $\pi_k$ can be resampled among greedy actions at every step. Thus the algorithm might break a tie by selecting this bad action during its waiting period. 
To prevent this, we introduce a global tie-breaking event that ensures no bad action is selected while its gap is zero:
\begin{align}
\label{eq: def global tie event}
\mathcal E^t_{t_n} \coloneqq
% \bigcap_{k=t_n}^{(T_n-1) \wedge t_{n+1}}
\bigcap_{k=t_n}^{\min\{ T_n-1 , t_{n+1} \}}
\left\{
\forall \, s \in \mathcal S : 
(s, \pi_k(s)) \notin \mathcal B_{\pi_{t_n}}
\right\} .
\end{align}

The joint event $\mathcal E^t_{t_n} \cap \mathcal E^s_{t_n} \cap \mathcal E^g_{t_n} \cap \mathcal E^l_{t_n}$ then guarantees that the greedy policy at time $t_{n+1}$ does not correspond to any pair in $\mathcal B_{\pi_{t_n}}$.


\begin{lemma}
\label{thm: event that yields positive gaps}
Consider the setting of \autoref{thm: fixed prob that bad actions dont become greedy}.
For every index $n \in \mathbb Z_{\ge 0}$ such that $t_n \ge K_N$,
the joint event $\mathcal E^t_{t_n} \cap \mathcal E^s_{t_n} \cap \mathcal E^g_{t_n} \cap \mathcal E^l_{t_n}$ guarantees that
% the greedy policy $\pi_k$ never coincides with a bad action, 
% meaning
for every time $k \in \{t_n, \dots, t_{n+1}\}$
and every state $s \in \mathcal S$,
\[
(s,\pi_{k}(s)) \notin \mathcal B_{\pi_{t_n}}
\]
\end{lemma}
\vspace{-0.6cm}

\begin{proof}
Fix an $n \in \mathbb Z_{\ge 0}$ such that $t_n \ge K_N$, and let $\pi \equiv \pi_{t_n}$.
For each state-action pair $(s,a) \in \mathcal B_\pi$, define the sequence $\big( d_k(s,a) \big)_{k=t_n}^{t_{n+1}}$ as in \autoref{def: gap variables}.
We track the evolution of a specific gap $d_k(s_\ell, a_\ell)$ across three phases.

By definition, all gaps are non-negative at time $t_n$ since $\pi_{t_n}$ is greedy with respect to $q_{t_n}$.
Moreover, on the joint event \(\mathcal E^t_{t_n} \cap \mathcal E^s_{t_n}\), the gap \(d_k(s_\ell,a_\ell)\) cannot decrease due to an earlier seed block \(\mathcal E^s_{b_j}(s_j,a_j)\) acting on a pair \(j<\ell\).
More precisely, let $(s_j, a)$ be the state-action pair updated by $\mathcal E^s_{b_j}(s_j,a_j)$. Since $j \neq \ell$ and $a_\ell \neq \pi(s_\ell)$, the case choice in \eqref{eq: def good prefix} implies that $(s_j,a) \neq (s_\ell,a_\ell)$.
So if $s_j \neq s_\ell$, then \(d_k(s_\ell,a_\ell)\) is unchanged;
if $s_j = s_\ell$ but $a \neq \pi(s_\ell)$, then \(d_k(s_\ell,a_\ell)\) is unchanged;
and if $s_j = s_\ell$ and $a = \pi(s_\ell)$, then \(d_k(s_\ell,a_\ell)\) increases exactly as in \eqref{eq: final lower bound case 1} and \eqref{eq: increase in d^b case B}.
Thus, by time \(b_\ell\),
either the gap associated with $(s_\ell,a_\ell)$ already hit \(\delta\), in which case \(\mathcal E^l_{t_n}(s_\ell,a_\ell)\) keeps it strictly positive up to time $t_{n+1}$ (\autoref{thm: event noise remains small}), or else \(0 \le d_{b_\ell}(s_\ell, a_\ell) < \delta\).

In the latter case, \autoref{thm: event to reach a_k margin} applied at time \(b_\ell\) on \(\mathcal E^s_{b_\ell}(s_\ell,a_\ell)\) yields
\[
d_{k}(s_\ell,a_\ell) > c_s A(b_\ell, k) > 0
\]
for every $k$ that satisfies 
% $b_\ell < k \le \big( \tau^s_{b_\ell}(s_\ell,a_\ell) \wedge (b_\ell + M) \wedge t_{n+1} \big)$.
$b_\ell < k \le \min \{ \tau^s_{b_\ell}(s_\ell,a_\ell) , b_\ell + M , t_{n+1} \}$.
If this gap does not reach \(\delta\) during its $M$-step seed block, then
\[
d_{b_\ell+M}(s_\ell,a_\ell) 
> c_s A(b_\ell, b_\ell+M)
\ge c_s M \alpha_{b_\ell+M}
\ge c_s M \alpha_{T_n}
\ge c_s \rho_\alpha M \alpha_{T_n}
\]
by the non-increasing property of $(\alpha_k)_{k \ge 0}$, and the fact that $\rho_\alpha \in (0,1]$.
Since later seed blocks can also not decrease this gap,
it either reaches $\delta$ before time $T_n$ and then \(\mathcal E^l_{t_n}(s_\ell,a_\ell)\) keeps it strictly positive up to $t_{n+1}$,
or else
\[
d_{T_n}(s_\ell,a_\ell)
> c_s \rho_\alpha M \alpha_{T_n}.
\]

Suppose the gap has not yet hit \(\delta\) by time \(T_n\). 
Since $T_n > t_n \ge K_N$,
\autoref{thm: from a_k margin to delta} guarantees that on the joint event
\[
\mathcal E^g_{T_n}(s_\ell,a_\ell) \cap
\big\{ c_s \rho_\alpha M \alpha_{T_n} \le d_{T_n}(s_\ell,a_\ell) < \delta \big\}
\]
that gap remains strictly positive until it reaches $\delta$ by time $T_n + N(T_n)$, or at least until the next transition time $t_{n+1}$, whichever happens first.
Again, if the gap reaches $\delta$ before $t_{n+1}$, then $\mathcal E^l_{t_n}(s_\ell, a_\ell)$ keeps it strictly positive up to $t_{n+1}$.

In summary,
for every $k \in \{t_n, \dots, T_n-1\}$,
the joint event $\mathcal E^t_{t_n} \cap \mathcal E^s_{t_n} \cap \mathcal E^l_{t_n}$
guarantees that $d_k(s,a) \ge 0$ for every $(s,a) \in \mathcal B_{\pi_{t_n}}$ and $(s,\pi_{k}(s)) \notin \mathcal B_{\pi_{t_n}}$ when a gap is zero in $s$.
As a result, if the transition occurs before time $T_n$, meaning if $t_{n+1} < T_n$, then, at time $t_{n+1}$,
for each pair $(s,a) \in \mathcal B_{\pi_{t_n}}$, 
either action $a$ is not greedy with respect to $q_{t_{n+1}}$,
or it is tied with $\pi_{t_{n}}(s)$ and $\mathcal E^t_{t_n}$ ensures the tie is broken in favour of an action outside $\mathcal B_{\pi_{t_n}}$, for instance $\pi_{t_{n}}(s)$.
In case $t_{n+1} \ge T_n$,
the joint event $\mathcal E^g_{t_n} \cap \mathcal E^l_{t_n}$
ensures that all gaps remain strictly positive
for the remaining time steps $k \in \{T_n, \dots, t_{n+1}\}$, and hence for every pair $(s,a) \in \mathcal B_{\pi_{t_n}}$,
\[
q_k(s, \pi_k) \ge q_k(s, \pi_{t_n}) > q_k(s,a).
\]
So no bad action can be greedy with respect to $q_k$ after time $T_n$.
Overall, the joint event $\mathcal E^t_{t_n} \cap \mathcal E^s_{t_n} \cap \mathcal E^g_{t_n} \cap \mathcal E^l_{t_n}$ guarantees that $(s,\pi_{k}(s)) \notin \mathcal B_{\pi_{t_n}}$ for all $k \in \{t_n, \dots, t_{n+1}\}$.
\end{proof}

 
%%%%%%%%%%%%%%%%%%%%%%%%%%%%%%%%%%%%%%%%%%%%%%%%%%%%%%%%
\subsubsection{Proof of \autoref{thm: fixed prob that bad actions dont become greedy}}

\begin{proof}
We first define deterministic probability bounds for the constituent events. Let
\begin{align}
p_g
:=
1-|\mathcal S\times\mathcal A|\,
\frac{8 c_\xi}{c^2\rho_\alpha^2}
\sum_{j=0}^{\infty}\frac{2^j}{(M+2^j)^2} .
\end{align}
Since
\[
\lim_{M \to \infty} \sum_{j = 0}^{\infty} \frac{2^j}{(M+2^j)^2} = 0
\]
we can choose $M \in \mathbb N$ sufficiently large such that $p_g \ge 1/2$.
Next, 
using the constant $\beta$ from the tie-breaking condition on the policy in \eqref{eq:greedy-pol-tie-breaking},
the minimum sampling probability $\sigma_{\min}$ from the strict exploring starts (\autoref{def: strict exploring starts}),
the minimum target gap $d_{\min}$ from \eqref{eq: min target gap},
and the upper bounds $c_v$ and $c_\xi$ on the variances of the perturbations $v_k$ and $\xi_k$,
define
\begin{align}
p_t &\coloneqq \beta^{|\mathcal S| } \in (0,1) ,
\\
p_s &:= \sigma_{\min} \frac{d_{\min}^2}{d_{\min}^2 + 16 c_v} \in (0,1) ,
\\
p_{\overline l} &\coloneqq 
|\mathcal S\times\mathcal A|
\frac{64 c_\xi}{\delta^2}\sum_{i=t_n}^\infty \alpha_i^2 .
\end{align}
Since the Robbins-Monro conditions require that \(\sum \alpha_i^2<\infty\), there exists a \(K_l<\infty\) such that for all \(t_n \ge K_l\), $p_{\overline l} \le \frac14 ( p_t p_s ) ^{|\mathcal S\times\mathcal A| M}$.


For every MC-O-PI solution there exists an almost surely finite time $K_\xi$ that satisfies \autoref{thm: properties noise gap}.
Fix an index $n \in \mathbb Z_{\ge 0}$ such that $t_n \ge \max\{K_\xi, K_N, K_l\}$.
Let $L := |\mathcal B_{\pi_{t_n}}|$ and assume that $L>0$, otherwise the result is trivial because \autoref{thm: strict improvement theorem} indicates that any transition would improve the policy.


On the joint event $\mathcal E^t_{t_n} \cap \mathcal E^s_{t_n} \cap \mathcal E^g_{t_n} \cap \mathcal E^l_{t_n}$, \autoref{thm: event that yields positive gaps} guarantees that no bad action becomes greedy. 
In other words, when $\pi_{t_n}$ is suboptimal, we obtain at some almost surely finite time $t_{n+1}$:
\begin{align}
\label{eq: condition that policy improves}
q_{\pi_{t_n}} ( s, \pi_{t_{n+1}} ) \ge q_{\pi_{t_n}} (s,\pi_{t_n} )
\end{align}
for every state $s\in S$.
By definition of the transition times, $q_{\pi_{t_n}} \neq q_{\pi_{t_{n+1}}}$.
Thus \autoref{thm: strict improvement theorem} proves that policy $\pi_{t_{n+1}}$ must be better than $\pi_{t_{n}}$, meaning $\pi_{t_{n+1}} \succ \pi_{t_n}$.
As a result, if $\pi_{t_n} \notin \mathcal P_*$, then
\[
\{\pi_{t_{n+1}} \succ \pi_{t_n} \}
\ \supseteq \
\mathcal E^t_{t_n} \cap \mathcal E^s_{t_n} \cap \mathcal E^g_{t_n} \cap \mathcal E^l_{t_n} .
\]
So it follows that, on the event $\{\pi_{t_n} \notin \mathcal P_*\}$,
\begin{align*}
\mathbb P\!\left(
\pi_{t_{n+1}} \succ \pi_{t_n}
\mid \mathcal F_{t_n}, \pi_{t_n}
\right)
\ge
\mathbb P ( \mathcal E^t_{t_n} \cap \mathcal E^s_{t_n} \cap \mathcal E^g_{t_n} \cap \mathcal E^l_{t_n} \mid \mathcal F_{t_n}, \pi_{t_n} ).
\end{align*}

Since $\mathcal E^t_{t_n} \cap \mathcal E^s_{t_n} \in \mathcal F_{T_n}$, we can decompose the probability as
\begin{align}
\nonumber
&\mathbb P \bigl(\mathcal E^t_{t_n} \cap \mathcal E^s_{t_n} \cap \mathcal E^g_{t_n} \cap \mathcal E^l_{t_n} \mid \mathcal F_{t_n} , \pi_{t_n} \bigr)
\\
\nonumber
&\quad \ge
\mathbb P \bigl(\mathcal E^t_{t_n} \cap \mathcal E^s_{t_n} \cap \mathcal E^g_{t_n} \mid \mathcal F_{t_n} , \pi_{t_n} \bigr) 
- \mathbb  P\bigl(\overline{\mathcal E^l_{t_n}} \mid \mathcal F_{t_n} , \pi_{t_n} \bigr)
\\
\label{eq: partial decomposition of the joint probabilities}
&\quad = \mathbb E \left[ \mathbf 1_{\mathcal E^t_{t_n} \cap \mathcal E^s_{t_n}} \mathbb P(\mathcal E^g_{t_n} \mid \mathcal F_{T_n}) \mid \mathcal F_{t_n}, \pi_{t_n} \right] - \mathbb P\bigl(\overline{\mathcal E^l_{t_n}} \mid \mathcal F_{t_n} , \pi_{t_n} \bigr)
\end{align}
Given that $T_n > t_n \ge \max \{K_\xi, K_N\}$, we obtain from \autoref{thm: from a_k margin to delta} and our specific choice of $M$ that
\begin{align*}
\mathbb P(\mathcal E^g_{t_n} \mid \mathcal F_{T_n}) 
&\ge 1 - \sum_{(s,a) \in \mathcal B_{\pi_{t_n}}} \mathbb P \bigl( \overline{\mathcal E^g_{T_n}(s,a)} \mid \mathcal F_{T_n} \bigr)
\\
&\ge 1 - |\mathcal B_{\pi_{t_n}}| \frac{8 c_\xi}{c^2 \rho_\alpha^2} \sum_{j = 0}^{\infty} \frac{2^j}{(M+ 2^j)^2}
\\
&\ge 1 - |\sset \times \aset| \frac{8 c_\xi}{c^2 \rho_\alpha^2} \sum_{j = 0}^{\infty} \frac{2^j}{(M+ 2^j)^2}
\\
&= p_g \ge \frac12 .
\end{align*}
Further, given that $t_n \ge \max \{K_\xi, K_l\}$, we obtain from \autoref{thm: event noise remains small} that
\begin{align*}
\mathbb P \bigl( \overline{\mathcal E^l_{t_n}} \mid \mathcal F_{t_n} , \pi_{t_n} \bigr)
&\le
\sum_{(s,a) \in \mathcal B_{\pi_{t_n}}} \mathbb P \bigl(\overline{\mathcal E^l_{t_n}(s,a)} \mid \mathcal F_{t_n}, \pi_{t_n} \bigr)
\\
&\le
|\mathcal B_{\pi_{t_n}}|
\frac{64 c_\xi}{\delta^2}\sum_{i={t_n}}^\infty \alpha_i^2
\\
&\le
|\mathcal S\times\mathcal A|
\frac{64 c_\xi}{\delta^2}\sum_{i={t_n}}^\infty \alpha_i^2
\\
&=
p_{\overline l}
\le
\frac14 ( p_t p_s)^{|\mathcal S \times \mathcal A| M} .
\end{align*}
Thus we can further lower-bound \eqref{eq: partial decomposition of the joint probabilities} by
\begin{align}
\label{eq: probability lower bound intermediary}
\mathbb  P(\mathcal E^t_{t_n} \cap \mathcal E^s_{t_n} \cap \mathcal E^g_{t_n} \cap \mathcal E^l_{t_n} \mid \mathcal F_{t_n}, \pi_{t_n} )
&\ge
p_g \, \mathbb P(\mathcal E^t_{t_n} \cap \mathcal E^s_{t_n} \mid \mathcal F_{t_n}, \pi_{t_n}) - p_{\overline l} .
\end{align}


It remains to lower-bound $\mathbb P(\mathcal E^t_{t_n} \cap \mathcal E^s_{t_n} \mid \mathcal F_{t_n}, \pi_{t_n})$. We do this by repeated conditioning in chronological order over the interval $\{t_n, \dots, T_n-1\}$.
Concretely, for $k \in \{t_n, \dots, T_n-1\}$, define
\[
\ell(k) := 1 + \left\lfloor \frac{k-t_n}{M} \right\rfloor \in \{1,\dots,L\} ,
\]
so that time \(k\) belongs to the \(M\)-step seed block associated with the state-action pair \((s_{\ell(k)},a_{\ell(k)})\) and starts at time
\[
b_{\ell(k)} = t_n + (\ell(k)-1) M.
\]
Using $\ell(k)$, we can define the one-step subevents 
of \(\mathcal E^s_{t_n}\) and \(\mathcal E^t_{t_n}\) that occur at each time $k \in \{t_n, \dots, T_n-1\}$.
Namely, if $q_{\pi_{t_n}} (s_{\ell(k)},\pi_{t_n} ) - q_{b_{\ell(k)}}(s_{\ell(k)},\pi_{t_n} ) \ge d_{\min}/2$,
then
\[
\mathcal E^s_{t_n}(k)
:=
\left\{
u_{k}(s_{\ell(k)},\pi_{t_n} )=1,\ 
v_{k}(s_{\ell(k)},\pi_{t_n} ) \ge -\dfrac{d_{\min}}{4}
\right\} ,
\]
otherwise
\[
\mathcal E^s_{t_n}(k)
:=
\left\{
u_{k}(s_{\ell(k)},a_{\ell(k)})=1,\ 
v_{k}(s_{\ell(k)},a_{\ell(k)}) \le \dfrac{d_{\min}}{4}
\right\} .
\]
Since \(k \ge b_{\ell(k)} \ge t_n\), the case selection is measurable given $\mathcal F_{k}$ and $\pi_{t_n}$.
Also, if $k \le t_{n+1}$, then
\[
\mathcal E^t_{t_n}(k)
:=
\left\{
\forall \, s \in \mathcal S : 
(s,\pi_k) \notin \mathcal B_{\pi_{t_n}}
\right\} .
\]
Otherwise it is the certain event
\[
\mathcal E^t_{t_n}(k)
:= \Omega ,
\]
where $\Omega$ is the set of all possible outcomes.
These one-step subevents provide the decompositions
\begin{align*}  
\mathcal E^s_{t_n} = \bigcap_{k=t_n}^{T_n-1}\mathcal E^s_{t_n}(k),
\\
\mathcal E^t_{t_n} = \bigcap_{k=t_n}^{T_n-1}\mathcal E^t_{t_n}(k).
\end{align*}


Now fix a time $k \in \{t_n, \dots, T_n\}$ and consider the joint event
\begin{align}
\label{eq: joint intermediate event}
\bigcap_{i=t_n}^{k-1} \bigl( \mathcal E^t_{t_n}(i)\cap \mathcal E^s_{t_n}(i) \bigr).
\end{align}
Regarding the tie-breaking event, we distinguish two cases.
First, suppose that $k \le t_{n+1}$. 
Every gap associated with a pair in \(\mathcal B_{\pi_{t_n}}\) is non-negative at time \(k\) because all gaps are non-negative at time \(t_n\) and seed blocks cannot decrease any gap below zero.
Therefore, if a pair \((s,a)\in\mathcal B_{\pi_{t_n}}\) is greedy at time \(k\), then necessarily
\[
d_k(s,a)=q_k(s,\pi_{t_n})-q_k(s,a)=0.
\]
Hence, \(\pi_{t_n}(s)\) is also greedy in state \(s\). It follows that, on the joint event \eqref{eq: joint intermediate event}, the tie-breaking condition \(\mathcal E^t_{t_n}(k)\) is satisfied whenever the policy at time \(k\) follows action \(\pi_{t_n}(s)\) in every state \(s\) where a bad action is greedy. By the action-selection condition in \eqref{eq:greedy-pol-tie-breaking} and the conditional independence of tie-breaking across states, we obtain
\[
\mathbb P\!\left(\mathcal E^t_{t_n}(k)\mid \mathcal F_k,\pi_{t_n}\right)
\ge
\beta^{|\mathcal S|}
=
p_t .
\]
Second, if \(k>t_{n+1}\), then $\mathcal E^t_{t_n}(k)$ is the certain event. So we also have
\[
\mathbb P\!\left(\mathcal E^t_{t_n}(k)\mid \mathcal F_k,\pi_{t_n}\right)=1\ge p_t .
\]
% In both cases,
% \[
% \mathbb P\!\left(\mathcal E^t_{t_n}(k)\mid \mathcal F_k,\pi_{t_n}\right)\ge p_t .
% \]
Regarding the seed event, the strict exploring starts conditions and inequalities \eqref{eq:noise-lower} and \eqref{eq:noise-upper} yield
\[
\mathbb P\!\left(\mathcal E^s_{t_n}(k)\mid \mathcal F_{k}, \pi_{t_n} \right)
\ge
\sigma_{\min}\frac{d_{\min}^2}{d_{\min}^2+16c_v}
=
p_s .
\]
Thus by using the tower property, we obtain
\begin{align*}
&\mathbb P\!\left(\mathcal E^t_{t_n}(k)\cap \mathcal E^s_{t_n}(k)\mid \mathcal F_{k}, \pi_{t_n }\right)
\\
&\qquad =
\mathbb E\!\left[ \mathbb E\!\left[ \mathbf 1_{\mathcal E^t_{t_n}(k)} \mathbf 1_{\mathcal E^s_{t_n}(k)} \mid \mathcal F_{k}, \pi_{t_n} \right] \middle| \mathcal F_k, \pi_{t_n}\right]
\\
&\qquad =
\mathbb E\!\left[ \mathbf 1_{\mathcal E^t_{t_n}(k)} \mathbb P \big( \mathcal E^s_{t_n}(k) \mid \mathcal F_{k}, \pi_{t_n} \big) \middle| \mathcal F_k, \pi_{t_n}\right]
\\
&\qquad \ge
p_s
\mathbb P\!\left(\mathcal E^t_{t_n}(k) \mid \mathcal F_{k}, \pi_{t_n} \right)
\ge
p_t p_s
\end{align*}
on $\bigcap_{i=t_n}^{k-1} \big( \mathcal E^t_{t_n}(i) \cap \mathcal E^s_{t_n}(i) \big)$.
Repeated conditioning in chronological order then yields
\begin{align*}
&\mathbb P\!\left(
\bigcap_{i=t_n}^{k} \bigl( \mathcal E^t_{t_n}(i) \cap \mathcal E^s_{t_n}(i) \bigr)
\ \middle|\ \mathcal F_{t_n}, \pi_{t_n}
\right)
\\
&\qquad=
\mathbb E\!\left[
\mathbf 1_{\cap_{i=t_n}^{k-1}(\mathcal E^t_{t_n}(i) \cap \mathcal E^s_{t_n}(i))}
\,
\mathbb P\!\left(\mathcal E^t_{t_n}(k)\cap \mathcal E^s_{t_n}(k) \mid \mathcal F_{k}, \pi_{t_n} \right)
\ \middle|\ \mathcal F_{t_n}, \pi_{t_n}
\right]
\\
&\qquad\ge
p_t p_s\,
\mathbb P\!\left(
\bigcap_{i=t_n}^{k-1} \bigl( \mathcal E^t_{t_n}(i) \cap \mathcal E^s_{t_n}(i) \bigr)
\ \middle|\ \mathcal F_{t_n}, \pi_{t_n}
\right).
\end{align*}
By iterating this inequality from time $t_n$ up to $T_n-1 = t_n + L M -1$,
and using $0 < p_t p_s \le 1$ along with $L \le |\mathcal S \times \aset|$,
we obtain 
\[
\mathbb P(\mathcal E^t_{t_n} \cap \mathcal E^s_{t_n} \mid \mathcal F_{t_n}, \pi_{t_n} )
\ge
(p_t p_s)^{LM}
\ge
(p_t p_s)^{|\mathcal S\times\mathcal A|\,M} .
\]

% Plugging this into \eqref{eq: probability lower bound intermediary},
Altogether, if we define a constant $p$ as
\begin{align}
\label{eq: def lower bound probability} 
p
:=
\frac14
(p_t \, p_s)^{|\mathcal S \times \mathcal A| \, M}
% \\
% \nonumber
=
\frac14\,
\left(
\beta^{|\mathcal S|}
\sigma_{\min}\frac{d_{\min}^2}{d_{\min}^2+16 c_v}
\right)^{|\mathcal S\times\mathcal A| \, M}
>0  ,
\end{align}
then for every MC-O-PI solution there exists an almost surely finite time $K \ge \max\{K_\xi, K_N, K_l\}$ such that for every $t_n \ge K$, the probability of strict improvement on the event $\{\pi_{t_n} \notin \mathcal P_*\}$ is
\begin{align*}
&\mathbb P\!\left(
\pi_{t_{n+1}} \succ \pi_{t_n}
\mid \mathcal F_{t_n}, \pi_{t_n}
\right)
\ge
p_g (p_t p_s)^{|\mathcal S\times\mathcal A|\,M} - p_{\overline l}
\ge
p .
\end{align*}
Since this lower bound is uniform over all realisations of $\pi_{t_n}$, the tower property gives
\begin{align*}
&\mathbb P\!\left(
\pi_{t_{n+1}} \succ \pi_{t_n}
\mid \mathcal F_{t_n}, \pi_{t_n} \notin \mathcal P_*
\right)
\\
&\qquad =
\mathbb E \! \left[ 
\mathbb P\!\left(
\pi_{t_{n+1}} \succ \pi_{t_n}
\mid \mathcal F_{t_n}, \pi_{t_n}
\right)
\middle| \mathcal F_{t_n}, \pi_{t_n} \notin \mathcal P_* \right]
\\
&\qquad \ge
\mathbb E [ p \mid \mathcal F_{t_n}, \pi_{t_n} \notin \mathcal P_* ]
\\
&\qquad =
p .
\end{align*}
\end{proof}

\autoref{thm: fixed prob that bad actions dont become greedy} established that, under uniform updates over all actions in each state, the probability of a policy improvement is eventually bounded below by a fixed constant. Since this lower bound is uniform over all realisations of the greedy policy, it remains valid when the current policy is optimal. In that case, however, no improvement is possible. Instead, the same joint event guarantees that no further transition occurs.

\vspace{0.3cm}

\begin{proposition}
\label{thm: prob policy remains optimal}
Consider the setting of \autoref{thm: fixed prob that bad actions dont become greedy}.
The constant $p>0$ and random time $K<\infty$ of \autoref{thm: fixed prob that bad actions dont become greedy} ensure that
for every index $n \in \mathbb Z_{\ge 0}$ such that $t_n \ge K$,
\[
\mathbb P (t_{n+1} = \infty \mid \mathcal F_{t_n}, \pi_{t_n} \in \mathcal P_* ) \ge p .
\]
\end{proposition}
\vspace{-0.6cm}
\begin{proof}
Fix an $n \in \mathbb Z_{\ge 0}$ such that $t_n \ge K$,
and assume that $\pi_{t_n}$ is optimal.
By definition of optimality, 
$q_{\pi_{t_n}}(s, \pi_{t_n} ) = \max_a q_{\pi_{t_n}}(s,a)$
for every state $s \in \mathcal S$.
As a result, for every pair
$(s,a) \notin \mathcal B_{\pi_{t_n}}$,
\begin{align}
\label{eq: implication of reached otpimal policy}
q_{\pi_{t_n}}(s,a) = q_{\pi_{t_n}}(s, \pi_{t_n}) .
\end{align}
Now suppose for contradiction that $t_{n+1} < \infty$
on the joint event $\mathcal E^t_{t_n} \cap \mathcal E^s_{t_n} \cap \mathcal E^g_{t_n} \cap \mathcal E^l_{t_n}$, as defined in \eqref{eq: def global seed event}--\eqref{eq: def global tie event}.
\autoref{thm: event that yields positive gaps} establishes that $(s,\pi_{t_{n+1}}(s)) \notin \mathcal B_{\pi_{t_n}}$ for every $s \in \mathcal S$. Hence, as in \eqref{eq: implication of reached otpimal policy},
\[
q_{\pi_{t_n}}(s,\pi_{t_{n+1}}) = q_{\pi_{t_n}}(s,\pi_{t_n}) .
\]
By \autoref{def: transition times}, $q_{\pi_{t_{n+1}}} \neq q_{\pi_{t_n}}$. Thus \autoref{thm: strict improvement theorem} implies that policy $\pi_{t_{n+1}}$ is better than $\pi_{t_{n}}$, which was already optimal. Given this is impossible, $t_{n+1}$ must be almost surely infinite on the joint event.
Consequently, as in the proof of \autoref{thm: fixed prob that bad actions dont become greedy},
\[
\mathbb P (t_{n+1} = \infty \mid \mathcal F_{t_n}, \pi_{t_n} \in \mathcal P_* )
\ge
\mathbb P(\mathcal E^t_{t_n} \cap \mathcal E^s_{t_n} \cap \mathcal E^g_{t_n} \cap \mathcal E^l_{t_n} \mid \mathcal F_{t_n}, \pi_{t_n} \in \mathcal P_* )
\ge p
\]
with $p$ defined in \eqref{eq: def lower bound probability}.
\end{proof}



Altogether, the proofs of \autoref{thm: fixed prob that bad actions dont become greedy} and \autoref{thm: prob policy remains optimal}
rely on three successive events.
These events ensure that stochastic perturbations cannot consistently
undermine the policy improvement driven by the mean-field dynamics
by forcing ``bad'' actions to become greedy.

% use advantage instead of gap?

The lock-in phenomenon is a direct application of \citep[Theorem~2.1]{Borkar2002OnApproximation}:
if the gap variable for a ``bad'' pair reaches a fixed positive margin infinitely often, 
it must eventually remain strictly positive.
However, since decreasing learning rates cause these gaps to be increasingly smaller at initialization,
we introduced a growth event to ensure they still reach this fixed margin.
This event extends the lock-in phenomenon to hold under a milder condition: gaps only need to reach a decreasing margin of size $O(\alpha_k)$ infinitely often.
Still, it requires an additional comparability condition on the learning rates so that, with a fixed probability, once a gap hits the $O(\alpha_k)$ margin, the remaining learning rates are sufficiently large to bring it up to the original fixed margin.
Finally, to guarantee gaps reach the decreasing $O(\alpha_k)$ margin within a fixed horizon, we used a seed event.
The length of this horizon determines the probability of the growth event.
By fixing a sufficiently long horizon, we can drive the probability of the growth event arbitrarily close to one, regardless of the unknown parameters in \eqref{eq: prob from a_k reach delta}.
The only trade-off for this extended horizon is a reduced probability of the seed event.
Yet, since that probability is fixed once the horizon is fixed, the overall probability of improvement remains fixed.


It is crucial for the seed event that updates use the initial-visit method to guarantee that all gaps initially increase.
% only one action be updated in any state at every step so 
% and that potential correlation in the episode noise does not break \eqref{eq: increase in d^b case B}.
% The structure of the stochastic perturbations $w_k$ is crucial for the seed event.
% Since the stochastic perturbations caused by episode simulations can be adversarially correlated, meaning they could 
% (which is not really an issue because uniform updates require initial-visit anyway in practice)
% to exactly cancel out the update in all but one state-action pair at every step (initial-visit).
However, uniform update distributions are irrelevant for this part of the proof because the strict exploring starts guarantee that we can always build the seed event.
On the contrary, the growth and lock-in events essentially enforce that the MC-O-PI solution behaves well because it does not wander away too much from the mean-field solution.
Thus they rely strongly on the uniformity of the update distributions that ensures the mean-field solution behaves well in the first place.
The update function is, however, irrelevant.
All that is necessary is that the stochastic perturbation $w_k$ is unbiased and that its variance is eventually bounded.



% Ultimately, the proof thus goes through the same
% under 
% initial-visit updates (allowing us to define the seed event),
% Robbins-Monro learning rates $(\alpha_k)$ that satisfy the comparability property,
% lower- and upper-bounded functions $(\mu_k)$ that are uniform over the actions in each state but do not necessarily integrate to one,
% and unbiased perturbations $(w_k)$ with eventually bounded variance.





%%%%%%%%%%%%%%%%%%%%%%%%%%%%%%%%%%%%%%%%%
%%%%%%%%%%%%%%%%%%%%%%%%%%%%%%%%%%%%%%%%%
\subsection{Global convergence of MC-O-PI}

\autoref{thm: fixed prob that bad actions dont become greedy} lower bounded the eventual probability of policy improvement,
while \autoref{thm: prob policy remains optimal} lower bounded the eventual probability of locking into the optimal region.
We now assemble these probabilistic results to prove that MC-O-PI converges to $\qopt$.

% say "lock-in"

\begin{proposition}
\label{thm: convergence mcopi uniform}
    Solutions of MC-O-PI (\autoref{def: mcopi}) converge almost surely to the optimal action values $\qopt$
    if
    the learning rates $(\alpha_k)_{k \ge 0}$ satisfy the Robbins-Monro conditions (\autoref{asp: stochastic approximation step size}) and the comparability assumption (\autoref{asp: additional conditions on learning rate}),
    the sampling distributions $(\sigma_k)_{k \ge 0}$ satisfy strict exploring starts (\autoref{def: strict exploring starts}),
    the update functions $(f_k)_{k \ge 0}$ are initial-visit (\autoref{def: initial visit updates}),
    and the update distributions $(\mu_k)_{k \ge 0}$ are uniform over all actions in each state (\autoref{asp: unbiased over policies}).
\end{proposition}
\vspace{-0.8cm}
\begin{proof}
Let $(q_k)_{k \ge 0}$ be a solution of MC-O-PI with associated policies $(\pi_k)_{k \ge 0}$ and transition times $(t_n)_{n \ge 0}$.
% We first show that there is an almost surely finite transition time $t_{n_*}$ such that $q_{\pi_k} = q_{\pi_*}$ for all $k \ge t_{n_*}$. We then use a standard stochastic approximation result to conclude that $q_k$ converges almost surely to $q_{\pi_*}$.
% 
By \autoref{thm: fixed prob that bad actions dont become greedy} and \autoref{thm: prob policy remains optimal}
there exists an almost surely finite time $K$ such that for every transition time $t_n \ge K$,
\[
\mathbb P (
\pi_{t_{n+1}} \succ \pi_{t_n}
\mid
\mathcal F_{t_n}, \pi_{t_n} \notin \mathcal P_*
)
\ge p ,
\]
and
\[
\mathbb P\!\left(
t_{n+1} = \infty
\mid
\mathcal F_{t_n}, \pi_{t_n} \in \mathcal P_*
\right)
\ge p ,
\]
where $p>0$ is the deterministic constant defined in \eqref{eq: def lower bound probability}.

Let $M \coloneqq |\mathcal{P}| < \infty$ denote the total number of deterministic policies.
Since any policy sequence can improve at most $M - 1$ times before reaching the optimal set $\mathcal{P}_*$, the probability of consecutive improvements until reaching $\mathcal{P}_*$ and staying there is eventually bounded below by $p^M$. In other words, for every $n$ such that $K \le t_n < \infty$,
\begin{align}
\label{eq:prob-lock-within-N-transitions}
\mathbb P\!\left(
t_{n+M} = \infty
\mid
\mathcal F_{t_n}
\right)
\ge p^M .
\end{align}

Fix an $n \in \mathbb Z_{\ge 0}$ such that $t_n \ge K$.
For every integer $i \ge 1$, on the event $\{t_{n+(i-1)M} < \infty\}$, 
equation \eqref{eq:prob-lock-within-N-transitions} applied at the transition index $n+(i-1)M$ yields
\begin{align*}
&\mathbb P (
t_{n+iM} < \infty
\mid
\mathcal F_{t_{n+(i-1)M}} )
\le 1-p^M .
\end{align*}
Taking conditional expectations with respect to $\mathcal{F}_{t_n}$ and repeatedly applying the tower property yields
\begin{align*}
\nonumber
&\mathbb P (
t_{n+iM} < \infty
\mid
\mathcal F_{t_{n}}
)
\\
&\qquad =
\mathbb E\!\left[
\mathbf 1_{\{t_{n+(i-1)M}<\infty\}}
\mathbb P (
t_{n+i M} < \infty
\mid
\mathcal F_{t_{n+(i-1)M}}
)
\middle|
\mathcal F_{t_{n}}
\right]
\\
&\qquad \le
(1-p^M)\,
\mathbb P (
t_{n+(i-1)M} < \infty
\mid
\mathcal F_{t_{n}}
)
\\
&\qquad \le
(1-p^M)^i
\end{align*}
Since $p > 0$, we have $1 - p^M<1$. Thus the right-hand side converges to 0 as $i$ grows to $\infty$. Consequently,
\[
\mathbb{P}\!\left( \bigcap_{i=1}^\infty \big\{ t_{n+iM} < \infty \big\} \;\middle|\; \mathcal{F}_{t_n} \right) = 0.
\]
Equivalently, with probability one, there exists a finite integer $i_*$ such that $t_{n+i_* M} = \infty$.

The event $\{t_{n+i_* M} = \infty\}$ guarantees that the sequence $(\pi_k)_{k \ge 0}$ eventually locks into the optimal set $\mathcal P_*$.
Therefore, there exists an index $n_*$ 
such that $\pi_{t_{n_*}} \in \mathcal P_*$ and $q_{\pi_k} = q_{\pi_*}$ for all $k \ge t_{n_*}$.
As a result, for every MC-O-PI solution $(q_k)_{k \ge 0}$ there exists an almost surely finite time $t_{n_*}$ after which
\begin{align}
    q_{k+1}
    =
    q_k + \alpha_k \big( \mu_k( q_{\pol_*} - q_k ) + w_k \big) .
\end{align}
Given that $\mu_k( q_{\pol_*} - q_k )$ is a descent direction on the potential $\|q_{\pol_*} - q_k\|_2^2$, 
that the learning rates satisfy the Robbins-Monro conditions,
that the sampling distributions satisfy strict exploring starts,
and that the stochastic perturbations are unbiased with bounded variance (\autoref{thm: properties of noise}),
Proposition~4.1 in \citep{Bertsekas1996Neuro-DynamicProgramming} guarantees that $q_k$ converges almost surely to $q_{\pi_*}$.
\end{proof}



% \begin{proposition}
\label{thm: convergence uniform}
    Solutions of MC-O-PI, as defined in \eqref{eq: complete difference inclusion}, converge to $q_{\best{\pol}}$ with probability one
    if the update distributions $(\update_k)_{k \geq 0}$ are uniform over all actions in each state (\autoref{asp: unbiased over policies})
    and the learning rates $(\lrate_k)_{k \geq 0}$ and update distributions $(\update_k)_{k \geq 0}$ satisfy the modified Robbins-Monro conditions (\autoref{asp: stochastic conditions on effective learning rate}).
\end{proposition}
% 
\begin{proof}
By \autoref{asp: standing}, $\sum_{k=0}^{\infty} \alpha_k^2 < \infty$ implies that $\lim_{k \to \infty} \alpha_k = 0$. Therefore, there exists an index $K$ such that $\alpha_k \in (0, 1]$ for all $k \ge K$.

For all $k \ge K$, \autoref{thm: improving policies} guarantees that the sequence of action values $(q_{\pi_k})$ is componentwise non-decreasing, given the uniform and strictly positive update distributions.
Since each $q_{\pi_k}$ belongs to the finite set $\{q_\pi : \pi \in \polset\}$, this sequence must stabilize. Thus, there exists a time step $K' \ge K$ and a policy $\pi$ such that $q_{\pi_k} = q_\pi$ for all $k \ge K'$.


As a result, for $k \ge K'$ the solution follows the linear system
\[
q_{k+1} = q_k + \alpha_k \mu(\sigma_k, \pi_k, f_k)( q_\pi - q_k)
\]
The Robbins-Monro and exploring starts conditions in \autoref{asp: standing} guarantee that the sequence $(q_k)$ converges to $q_\pi$.

Since $\lim_{k \to \infty} q_k = q_\pi$, \autoref{thm: asymptotic set of policies} establishes that there exists a time step $K'' \ge K'$ such that $\grpol(q_k) \subseteq \grpol(q_\pi)$ for all $k \ge K''$.
By definition, the active policy $\pi_k$ always satisfies $\pi_k \in \grpol(q_k)$. 
Therefore, for all $k \ge K''$
\[
\pi_k \in \grpol(q_\pi)
\]
Since $q_\pi = q_{\pi_k}$ for all $k \ge K'$, substituting this into the relation above yields $\pi_k \in \grpol(q_{\pi_k})$.
This confirms that for all $k \ge K''$, the policy $\pi_k$ is greedy with respect to its own action values, and hence satisfies the Bellman optimality equation $q_{\pi_k} = T_{\pi_k} q_{\pi_k} = T q_{\pi_k}$.
Consequently, $\pi_k$ is an optimal policy and the sequence $(q_k)$ converges to $q_\pi = \qopt$.
\end{proof}
% \begin{proof}
% Consecutive elements in every solution of exact MC-O-PI satisfy equation \eqref{eq: definition of q+ uniform section}
% and by assumption the update distributions $(\update_k)$ are uniform over all actions in each state.
% \autoref{thm: unbiased mces policy improves} thus guarantees that the value of consecutive policies is non-decreasing. Since this value is also upper bounded by $\qopt$, it converges.
% As a result, for every solution of exact MC-O-PI
% % regardless of the learning rates, start distributions and update functions,
% there exists a time step $K$ and a policy $\pol \in \polset$ such that, for all $k \geq K$, $q_k \in \graval(\pol)$ and the solution follows the system:
% \begin{align}
% \label{eq: equivalent system on the long term}
%     q_\kpo
%     = 
%     q_k + \lrate_k \update(\start_k, \pol, f_k) ( q_\pol - q_k ) . 
% \end{align}
% The modified Robbins-Monro conditions on $(\lrate_k)$ and $(\update_k)$ guarantee this system converges to $q_\pol$, which implies that $\qpol \in \graval(\pol)$.
% % As explained in \autoref{rmk: only qopt is inside its own region},
% Every policy that satisfies $\qpol \in \graval(\pol)$, satisfies the Bellman optimality equation.
% Thus $\pol = \best{\pol}$ and $\qopt$ is globally asymptotically stable for exact MC-O-PI, regardless of the first learning rate $\lrate_0$, start distribution $\start_0$ and update function $f_0$.
% Therefore, by \autoref{thm: if det converges then stoch converges}, solutions of MC-O-PI converge to $\qopt$ with probability one.
% \end{proof}


% As noted in \autoref{rmk: eff lr in the limit}, uniformity is only required on the effective learning rates in the limit. 
% However, we impose it on the update distributions in \autoref{thm: convergence uniform} for clarity. Specifically, verifying the condition on the sequence $(\update_k)$ is more evident than on the products $(\lrate_k \update_k)$.
% Besides, we assume $\lrate_k \in \real_{>0}$ unless specified otherwise. So ensuring uniformity of $(\lrate_k \update_k)$ ultimately reduces to verifying it on $(\update_k)$.



%%%%%%%%%%%%%%%%%%%%%%%%%%%%%%%%%%%%%%%%%
%%%%%%%%%%%%%%%%%%%%%%%%%%%%%%%%%%%%%%%%%
\section{Discussion}

\autoref{thm: convergence mcopi uniform} proved that initial-visit MC-O-PI converges almost surely to the optimal action values of a discounted or episodic MDP
when updates are uniform over the actions in each state
or, equivalently, when updates are uniform over the set of policies.
% provided
% that the learning rates satisfy the Robbins-Monro conditions and the comparability property,
% % the update functions are initial-visit,
% that the sampling distributions satisfy strict exploring starts,
% and that the resulting update distributions are uniform over all actions in each state.
% 
This section explores the implications of this result.
We first argue that some conditions can likely be relaxed because they are not essential for the underlying improvement process.
We then explain that the uniformity condition effectively requires using initial-visit updates, which does not necessarily slow down convergence compared to first-visit updates.
Finally, we show how to turn non-uniform updates into uniform updates in order to guarantee convergence.


%%%%%%%%%%%%%%%%%%%%%%%%%%%%%%%%%%%%%%%%%%%%%%%%%%%%%%%%%%%%%%%%%%%
%%%%%%%%%%%%%%%%%%%%%%%%%%%%%%%%%%%%%%%%%%%%%%%%%%%%%%%%%%%%%%%%%%%
\subsection{Relaxing conditions}
\label{sec: relaxing conditions}

% - assumption that limit set reduces to a point => this is assumption on solution, not on algorithm

% relax
% - policy samping => update only policy in state that changed, then dont even need event $\mathcal E^t_k$
% - essentially just need comparison prop of leanring rates to ensure reach fixed distance into a region, if can show other way, then just need Borkar lock in, i.e. \autoref{thm: event noise remains small}

% Technically, 
The proof of \autoref{thm: convergence mcopi uniform}
requires a uniform lower-bound on the distribution of the greedy policy at each step, strict exploring starts and non-increasing, comparable learning rates.
However, these conditions are artefacts of the method for the proof rather than necessary conditions for the underlying policy improvement process.


For instance, the global seed event \eqref{eq: def global seed event} updates only one state-action pair per step. Yet the greedy policy can be updated for all states at every step.
Sampling the greedy policy, as described in \eqref{eq:greedy-pol-tie-breaking}, 
therefore prevents
adversarial tie-breaking from always worsening the policy by choosing $\pi_{k+1}(s)=a$
whenever
% if after the $k^{\text{th}}$ update,
$q_{k+1}(s,\pi_k) = q_{k+1}(s,a)$ but $q_{\pi_k}(s,\pi_k) > q_{\pi_k}(s,a)$.
% the gap of a pair in $\mathcal B_\pi$ is zero after an update,
% , and thereby invalidate \autoref{thm: fixed prob that bad actions dont become greedy}.
% (see the discussion from equation \eqref{eq: joint intermediate event} onwards)
However, such ties are rare in practice.
% except if the algorithm is initialised in one.
Moreover, it is unrealistic to update the greedy policy for all states at every step.
A more practical method under initial-visit updates is to only recompute the policy in the state that was updated at step $k$,
by comparing $q_{k+1}(s, \pi_{k})$ with $q_{k+1}(s, a)$ where $u_{k}(s,a) = 1$.
This method simplifies the proof of \autoref{thm: convergence mcopi uniform} as the global tie-breaking event \eqref{eq: def global tie event} is not necessary anymore,
given that ties are broken once, right after the transition, 
and remain unchanged until the corresponding seed event makes the policy-worsening action non-greedy.
% until the pair is updated by the corresponding seed event,
% which can only increase the gap and hence break the tie in favour of the currently greedy action.


Further, the strict exploring starts, and the non-increasing and comparability conditions on the learning rates ensure that the stochastic perturbations cannot persistently offset the mean-field drift.
Specifically, they guarantee through the seed and growth events \eqref{eq: def global seed event}--\eqref{eq: def global grow event} that
for every policy $\pi \in \mathcal P$,
the stochastic perturbations cannot force the sequence
$(q_k(s,\pi) - q_k(s,a))_{k \ge 0}$ to converge to zero 
on the subsequence of times $\{k \ge 0 : \pi_k = \pi\}$
if $q_{\pi}(s,\pi) > q_{\pi}(s,a)$.
However, given that the accumulated stochastic perturbations form a converging martingale sequence, and hence $\lim_{K \to \infty} \sum_{k \ge K} \alpha_k w_k = 0$ (\autoref{thm: effects of noise decrease}), we expect that there exists another way to establish that the perturbations do not contain enough energy to consistently counteract the mean-field dynamics.
% , likely using an equivalent of the Robbins-Siegmund Theorem.

% If can prove another way, then can use the lock-in phenomenon and the combined stability-ODE method described in Section \ref{sec: limitations of existing methods}.


Since the monotone improvement that drives convergence when updates are uniform does not hinge on these conditions, we expect that they can be relaxed.
% Nevertheless, some constraints on the update distributions will always remain necessary as they effectively scale the learning rates in the MC-O-PI difference inclusion \eqref{eq: mcopi not yet sri}.
% leading to different learning rates for different state-action pairs.
% Namely, convergence to a fixed point requires Robbins–Monro–like conditions on the products $(\lrate_k \update_k)$, not on $(\lrate_k)$ alone.

% To see this, fix a policy $\pi$ and consider the system
% \begin{align}
% \label{eq: single system}
%     q_\kpo 
%     % &= q_k + \lrate_k \update_k (q_\pi - q_k + w_k)
%     % \\
%     &= q_k + \lrate_k \update_k (q_\pi - q_k) + \lrate_k w_k .
%     % + \lrate_k (v_k(\update_k + v'_k) + v'_k(q_\pi - q_k))
% \end{align}
% % where
% % \begin{align}
% %     w_k
% %     =
% %     u_k v_k + (u_k - \mu_k) ( q_{\pol} - q_k ) .
% % \end{align}
% Unrolling over $n$ steps yields
% \begin{multline*}
% % \label{eq: single system unrolled}
%     q_{k+n}
%     = \Bigg( \prod_{i=k}^{k+n-1} (1-\lrate_i \update_i) \Bigg) q_k 
%     + \Bigg( 1 - \prod_{i=k}^{k+n-1} (1-\lrate_i \update_i) \Bigg) q_\pi
%     \\
%     + \sum_{j=k}^{k+n-1} \prod_{i=j+1}^{k+n-1} (1-\lrate_i \update_i) \lrate_j w_j .
% \end{multline*}
% For $q_{k+n}$ to converge to $q_\pi$ as $n$ grows to $\infty$, three conditions must be satisfied.
% First, to ensure all state-action pairs are updated infinitely often,
% % Therefore, for every state-action pair $(s,a)$, the update distributions $(\update_k)$ must satisfy
% \begin{align}
% \label{eq: condition infinite visits}
%     \sum_{k=0}^{\infty} \update_k(s,a) = \infty
%     \qquad \forall \, s \in \sset, \forall \, a \in \aset(s) .
% \end{align}
% Second, to ensure the products $\prod_{i=k}^{k+n-1} (1-\lrate_i \update_i)$ decay to zero for increasing $n$,
% \begin{align}
% \label{eq: condition infinite effective learning rate}
%     \sum_{k=0}^{\infty} \lrate_k \update_k(s,a) = \infty
%     \qquad \forall \, s \in \sset, \forall \, a \in \aset(s) .
% \end{align}
% Finally, to ensure the accumulated noise vanishes,
% % This is ensured by the standard square-summability condition on the learning rates:
% \begin{align}
%     \sum_{k=0}^\infty \lrate_k^2 < \infty .
% \end{align}
% % If we imposed $\sum_k \lrate_k^2 \update_k^2 < \infty$ instead, we could choose $\lrate_k = \update_k(s,a) = 1/\sqrt{k}$ for each state-action pair, still satisfy conditions \eqref{eq: condition infinite visits} and \eqref{eq: condition infinite effective learning rate}. However, the noise term in \eqref{eq: single system unrolled} would become
% % \begin{align}
% %     \sum_{j=k}^{k+n-1} \prod_{i=j+1}^{k+n-1} (1-\frac{1}{i}) \frac{1}{\sqrt{j}} \xi_j
% %     &= \sum_{j=k}^{k+n-1} \frac{j}{k+n+1} \frac{1}{\sqrt{j}} \xi_j
% %     \\
% %     &= \frac{1}{k+n+1} \sum_{j=k}^{k+n-1} \sqrt{j} \xi_j
% %     \\
% %     &\eqqcolon N_{k,n}
% % \end{align}
% % We see that
% % \begin{align}
% %     \E[N_{k,n}^2] = \E \Bigg[ \Big( \frac{1}{k+n+1} \sum_{j=k}^{k+n-1} \sqrt{j} \xi_j \Big)^2 \Bigg]
% % \end{align}
% % Since the cross terms vanish in expectation, 
% % and as a result, $q_k$ would not converge to $q_\pi$
% % 
% % since $\sum_k \lrate_k^2 = \infty$, the noise would not be almost surely bounded. In other words, the noise alone would be able to pull solution arbitrarily far from $q_\pi$ infinitely often, thereby preventing convergence.
% % 
% Note that \eqref{eq: condition infinite visits} does not automatically imply \eqref{eq: condition infinite effective learning rate}.
% For instance, under \eqref{eq: condition infinite visits} and the Robbins-Monro conditions on the learning rates $(\lrate_k)$, we could choose $\lrate_k = \update_k(s,a) = 1/k$. System \eqref{eq: single system} would then converge prematurely due to Zeno-like behaviour because the products $\prod_{i=k}^{k+n-1} (1-\lrate_i \update_i)$ would not decay to 0.


% Altogether, the minimal requirements on the learning rates and the update distributions are
% % for every state $s \in \sset$ and action $a \in \aset(s)$
% \begin{align}
%     \label{eq: xtn robbins-monro first}
%     &\sum_{k =0}^\infty \lrate_k \update_k(s,a) = \infty 
%     \qquad \forall \, s \in \sset, \forall \, a \in \aset(s) ,
%     \\
%     \label{eq: xtn robbins-monro second}
%     &\sum_{k =0}^\infty \lrate_k^2 < \infty .
% \end{align}
% since \eqref{eq: xtn robbins-monro second} implies that, eventually, $\alpha_k \mu_k < \mu_k$, and \eqref{eq: xtn robbins-monro first} then lower bounds \eqref{eq: condition infinite visits}.
% We expect the results from Section~\ref{sec: results} to extend under these milder conditions.


% e.g. update distribution
% e.g. policy
% the GLIE conditions (\autoref{def: glie}) require that in the limit the policy $\pol$ is greedy with respect to the estimate $q$ so that solutions stop moving once they reach the optimal policy $\best{\pol}$ and action values $\qopt$.
% Since we study the asymptotic behaviour of MC-O-PI, we will assume that the policy is greedy at every step.

% Tsitsiklis notes that MC-O-PI's asymptotic behaviour depends on the asymptotic behaviour of the effective learning rates $(\lrate_k \update_k)$, not on either component alone.
% For instance, if the update distributions are non-uniform but constant, say $\update_k = \update$ for all $k$, and the learning rates take the form $\lrate_k(s,a) = 1/n_k(s,a)$, where $n_k(s,a)$ denotes the number of updates of $(s,a)$ in the first $k$ steps,
% then $\lrate_k \update_k$ converges to $1/k$ because for each state-action pair $n_k(s,a)$ tends to $k \update(s,a)$. As a result, \autoref{thm: c(q) goes to zero} still applies.

%%%%%%%%%%%%%%%%%%%%%%%%%%%%%%%%%%%%%%%%%%%%%%%%%%%%%%%%%%%%%%%%%%%
%%%%%%%%%%%%%%%%%%%%%%%%%%%%%%%%%%%%%%%%%%%%%%%%%%%%%%%%%%%%%%%%%%%
\subsection{Uniform updates effectively force initial-visits}


Using first- or every-visit updates causes the update distribution $\mu_k$ to depend on the transition dynamics induced by policy $\pi_k$ on the underlying MDP.
These dynamics are unknown when the MDP is itself unknown.
Therefore, requiring uniform update distributions effectively forces the use of initial-visit updates,
% because we cannot determine how the initial sampling distribution $\sigma_k$ will propagate through the MDP under policy $\pi_k$. 
because by updating only the initial state-action pair of each episode, we ensure $\mu_k$ equals $\sigma_k$ and thus retain control.

This control comes at the cost of data efficiency: the initial-visit method discards within-episode transitions that the first-visit method would reuse. However, more updates per episode do not automatically translate into faster policy convergence.
The extra updates induce a bias in the update distribution that can generate sliding modes in the action-value estimates $(q_k)$ and thereby slow down the convergence of the policy sequence $(\pi_k)$. Ultimately, it is the convergence speed of this policy sequence that truly matters for identifying the optimal behaviour of the MDP.


% describe MDP
To illustrate this, consider the MDP in \autoref{fig: experiment smallest MDP for iv vs fv} with a discount factor $\discount = 0.5$. Action \textit{left} yields a reward of one and \textit{right} yields a reward of zero.
% The best policy $\best{\pol}$ follows the \textit{left} action, whereas the worst policy $\worst{\pol}$ follows the \textit{right} action.
% describe learning algorithm
Under a uniform sampling distribution $\sigma_k$, the intial-visit method updates exactly one state-action pair per episode, whereas the first-visit method updates on average $1.5$ pairs.
With a learning rate $\alpha_k = 1/(k+100)$, the first-visit method also benefits from larger effective steps because it updates state-action pairs more frequently.
% Despite this numerical advantage, experiments show that IV can convergence to $\pi^*$ sooner.


Despite this apparent advantage, the initial-visit method is more stable than first-visit when MC-O-PI is initialised above $\qopt$.
Namely, initial-visit solutions follow straight lines that always point towards better policies.
In contrast, first-visit solutions exhibit a sliding mode that causes high-frequency oscillations in the policy sequence and traps $q_k$ near the boundary for longer.
As a result, first-visit can take longer to settle on the optimal policy.
For instance, when averaging over 1000 solutions starting in $q'$, first-visit ones require on average 43 updates to settle on $\pi_*$, whereas initial-visit ones always remain in $\region(\best{\pi})$.
When starting in $q$, the solutions take 68 and 70 updates to settle on $\pi_*$.

% \hrule
\begin{figure}[h!]
    \centering
    \begin{subfigure}[b]{0.3\textwidth}
        \centering
        \begin{tikzpicture}[
            scale=1, auto, node distance=2cm, 
            every edge/.style={draw, <-, font=\small},
            state/.style={circle, fill=mylightgrey, draw=black, text=black}
        ]
            % Nodes
            \node[state] (s1) at (0,0) {$s$};
            % Edges
            \path[->]
                (s1) edge[out=160, in=200, loop, looseness=8, draw=black] node[pos=0.5, left, text=black] {1} (s1)
                (s1) edge[out=20, in=340, loop, looseness=8, draw=black] node[pos=0.5, right, text=black] {0} (s1) ;
                % (s1) edge[out=145, in=215, loop, looseness=7] 
                %     node[pos=0.5, right] {$\best{a}$} 
                %     node[pos=0.5, left] {1} 
                % (s1)
                % (s1) edge[out=35, in=325, loop, looseness=7] 
                %     node[pos=0.5, left] {$\worst{a}$} 
                %     node[pos=0.5, right] {0} 
                % (s1);
        \end{tikzpicture}
        \caption{MDP}
        \label{fig: experiment smallest MDP}
    \end{subfigure}%
    \hfill
    \begin{subfigure}[b]{0.3\textwidth}
        \centering
        \begin{tikzpicture}[
            scale=1, auto, node distance=2cm, 
            every edge/.style={draw, <-, font=\small},
            state/.style={circle, fill=mylightgrey, draw=black, text=black}
        ]
            % Nodes
            \node[state] (s1) at (0,0) {$s$};
            % Edges
            \path[->]
                (s1) edge[out=160, in=200, loop, looseness=8, draw=black] node[pos=0.5, left, text=black] {1} (s1)
                (s1) edge[out=20, in=340, loop, looseness=8, draw=mylightgrey] node[pos=0.5, right, text=mylightgrey] {0} (s1) ;
                % (s1) edge[out=160, in=200, loop, looseness=8] node[pos=0.5, left] {$\begin{aligned}
                %     & \best{a} \\
                %     & 1
                % \end{aligned}$} (s1)
                % (s1) edge[out=20, in=340, loop, looseness=8, draw=mylightgrey] node[pos=0.5, right, text=mylightgrey] {$\begin{aligned}
                %     & \worst{a} \\
                %     & 0
                % \end{aligned}$} (s1);
        \end{tikzpicture}
        \caption{Policy $\best{\pol}$}
    \end{subfigure}%
    \hfill
    \begin{subfigure}[b]{0.3\textwidth}
        \centering
        \begin{tikzpicture}[
            scale=1, auto, node distance=2cm, 
            every edge/.style={draw, <-, font=\small},
            state/.style={circle, fill=mylightgrey, draw=black, text=black}
        ]
            % Nodes
            \node[state] (s1) at (0,0) {$s$};
            % Edges
            \path[->]
                (s1) edge[out=160, in=200, loop, looseness=8, draw=mylightgrey] node[pos=0.5, left, text=mylightgrey] {1} (s1)
                (s1) edge[out=20, in=340, loop, looseness=8, draw=black] node[pos=0.5, right, text=black] {0} (s1) ;
                % (s1) edge[out=160, in=200, loop, looseness=8, draw=mylightgrey] node[pos=0.5, left, text=mylightgrey] {$\begin{aligned}
                %     & \best{a} \\
                %     & 1
                % \end{aligned}$} (s1)
                % (s1) edge[out=20, in=340, loop, looseness=8] node[pos=0.5, right] {$\begin{aligned}
                %     & \worst{a} \\
                %     & 0
                % \end{aligned}$} (s1);
        \end{tikzpicture}
        \caption{Policy $\worst{\pol}$}
    \end{subfigure}%
    \caption{MDP with two policies.}
    \label{fig: single state mdp}
\end{figure}
\vspace{-0.4cm}
\hrulefill
\begin{figure}[h!]
    \centering
    \begin{subfigure}[b]{\textwidth}
        \centering
        \begin{tikzpicture}[
            scale=1, auto, node distance=2cm, 
            every edge/.style={draw, <-, font=\small},
            state/.style={circle, fill=mylightgrey, draw=black, text=black}
        ]
            % Nodes
            \node[state] (s1) at (0,0) {$s$};
            % Edges
            \path[->]
                (s1) edge[out=160, in=200, loop, looseness=8, draw=black] node[pos=0.5, left, text=black] {\textcolor{mygrey}{$\best{\pi}(s) :=$} 1} (s1)
                (s1) edge[out=20, in=340, loop, looseness=8, draw=black] node[pos=0.5, right, text=black] {0 \textcolor{mygrey}{$=: \worst{\pi}(s)$}} (s1) ;
                % (s1) edge[out=145, in=215, loop, looseness=7] 
                %     node[pos=0.5, right] {$\best{a}$} 
                %     node[pos=0.5, left] {1} 
                % (s1)
                % (s1) edge[out=35, in=325, loop, looseness=7] 
                %     node[pos=0.5, left] {$\worst{a}$} 
                %     node[pos=0.5, right] {0} 
                % (s1);
        \end{tikzpicture}
        \caption{MDP}
        \label{fig: experiment smallest MDP for iv vs fv}
    \end{subfigure}%
    \\
    \vspace{0.2cm}
    % --- Subfigure 1: Streamlines ---
    \begin{subfigure}[b]{0.48\textwidth}
        \begin{tikzpicture}[scale=0.9]

            \draw[->] (0, 0) -- (4.5, 0) node[right, font=\small] {$q(s,\pol_*)$};
            \draw[->] (0, 0) -- (0, 4.5) node[above, font=\small] {$q(s,\overline \pol)$};
            
            \begin{axis}[
                width=\linewidth,
                % title={Initial-visit solutions},
                % xlabel={$q(s, \pi_*(s))$},
                % ylabel={$q(s, \overline \pi(s))$},
                tick label style={font=\small},
                label style={font=\small},
                xmin=0, xmax=5,
                ymin=-0.5, ymax=4.5,
                % grid=off,
                scale=1.23,
                axis line style={draw=none},
                axis equal image,
                unbounded coords=jump, % This handles the NaN values in your CSV
                xtick={1, 2,4},
                % xticklabels={1, $\frac{1}{1-\gamma}$, 4},
                ytick={0, 1,4}
            ]

                \addplot [black, thin, dotted, domain=-0.5:5, samples=2] {x};
                
                % 1. Plot the streamlines from CSV
                \addplot[mygrey, thin, forget plot] table [col sep=comma, x index=0, y index=1] {figures/uniform/streamlines-iv.csv};

                % 2. Plot the simulation trajectory
                \addplot[CambridgeCoreBlue, thick] table [col sep=comma, x index=0, y index=1] {figures/uniform/simulation-iv-above.csv};

                \addplot[CambridgeCoreBlue, thick] table [col sep=comma, x index=0, y index=1] {figures/uniform/simulation-iv-below.csv};

                % 3. Mark Attractors (x1 and x2 based on your gamma=0.5)
                % \filldraw[fill=black, draw=black] (2, 1) circle (1.5pt) node[right, text=black, font=\small] {$q_{\pi_*}$};
                \node[label={right:$q_{\pi_*}$}, circle, fill, inner sep=1.5pt] at (axis cs:2, 1) {};
                
                \node[label={right:$q_{\overline \pi}$}, circle, fill, inner sep=1.5pt] at (axis cs:1,0) {};
                
                % % 4. Start Points with Labels
                % % Point A: Using your x0 = [3.75, 4]
                \node[label={left:$q$}, circle, fill, inner sep=1.5pt, CambridgeCoreBlue] at (axis cs:3.75, 4) {};
                
                % % Point B: Placeholder (Update coordinates as needed)
                \node[label={below:$q'$}, circle, fill, inner sep=1.5pt, CambridgeCoreBlue] at (axis cs:4, 3.75) {};
            \end{axis}
        \end{tikzpicture}
        \caption{Initial-visit solutions}
        \label{fig: iv versus fv simulations IV}
    \end{subfigure}
    \hfill
    % --- Subfigure 2: Placeholder or Second Simulation ---
    \begin{subfigure}[b]{0.48\textwidth}
        \begin{tikzpicture}[scale=0.9]

            \draw[->] (0, 0) -- (4.5, 0) node[right, font=\small] {$q(s,\pol_*)$};
            \draw[->] (0, 0) -- (0, 4.5) node[above, font=\small] {$q(s,\overline \pol)$};
            
            \begin{axis}[
                width=\linewidth,
                % title={First-visit solutions},
                % xlabel={$q(s, \pi_*(s))$},
                % ylabel={$q(s, \overline \pi(s))$},
                tick label style={font=\small},
                label style={font=\small},
                xmin=0, xmax=5,
                ymin=-0.5, ymax=4.5,
                % grid=off,
                scale=1.23,
                axis line style={draw=none},
                axis equal image,
                unbounded coords=jump, % This handles the NaN values in your CSV
                xtick={1, 2,4},
                % xticklabels={1, $\frac{1}{1-\gamma}$, 4},
                ytick={0, 1,4}
            ]

                \addplot [black, thin, dotted, domain=-0.5:5, samples=2] {x};
                
                % 1. Plot the streamlines from CSV
                \addplot[mygrey, thin, forget plot] table [col sep=comma, x index=0, y index=1] {figures/uniform/streamlines-fv.csv};

                % 2. Plot the simulation trajectory
                \addplot[CambridgeCoreBlue, thick] table [col sep=comma, x index=0, y index=1] {figures/uniform/simulation-fv-above.csv};

                \addplot[CambridgeCoreBlue, thick] table [col sep=comma, x index=0, y index=1] {figures/uniform/simulation-fv-below.csv};
                
                % 3. Mark Attractors (x1 and x2 based on your gamma=0.5)
                \node[label={right:$q_{\pi_*}$}, circle, fill, inner sep=1.5pt] at (axis cs:2, 1) {};
                
                \node[label={right:$q_{\overline \pi}$}, circle, fill, inner sep=1.5pt] at (axis cs:1,0) {};
                
                % 4. Start Points with Labels
                % Point A: Using your x0 = [3.75, 4]
                \node[label={left:$q$}, circle, fill, inner sep=1.5pt, CambridgeCoreBlue] at (axis cs:3.75, 4) {};
                
                % Point B: Placeholder (Update coordinates as needed)
                \node[label={below:$q'$}, circle, fill, inner sep=1.5pt, CambridgeCoreBlue] at (axis cs:4, 3.75) {};
            \end{axis}
        \end{tikzpicture}
        \caption{First-visit solutions}
        \label{fig: iv versus fv simulations FV}
    \end{subfigure}
    \caption{By avoiding sliding modes, initial-visit solutions sometimes converge to the optimal policy $\best{\pi}$ faster than first-visit solutions despite fewer updates and smaller effective learning rates.}
    \label{fig: iv versus fv simulations}
\end{figure}


% (\autoref{tab: time to converge to pol opt}).
% \textcolor{red}{remove table, integrate in text}
% \begin{table}[b]
%     \centering
%     \begin{tabular}{|r|cc|}
%     \hline
%         Start & IV solution & FV solution  \\
%         $A$ & $69 \pm 52$ & $68 \pm 40$ \\
%         $B$ & $0 \pm 0$ & $43 \pm 58$ \\
%         \hline
%     \end{tabular}
%     \caption{Average number of steps before $(\pi_k)$ converges to $\best{\pi}$ for 1000 simulations.}
%     \label{tab: time to converge to pol opt}
% \end{table}

Overall, this experiment illustrates that the initial-visit method can be practical in real-world applications. 
By avoiding sliding modes, it sometimes converges to the optimal policy $\best{\pi}$ as fast as, or faster than, the first-visit method despite fewer updates and smaller effective learning rates.
Besides, the initial-visit method is more reliable for early stopping. Namely, if training stops early, the last greedy policy of an initial-visit solution is typically the best policy encountered so far, whereas the first-visit solution may be oscillating across a decision boundary, making it unclear which policy is best.



% if we use learning rate $1/n(s,a)$
% actually helps FV because learning rate decreases the greedy bias

% Namely, since the policy sequence transitions almost smoothly instead of oscillating, terminating the algorithm prematurely yields a policy that reliably reflects the best policy found so far. In contrast, stopping a first-visit trajectory early risks catching it across a decision boundary, leaving it unclear which policy is actually best.


%%%%%%%%%%%%%%%%%%%%%%%%%%%%%%%%%%%%%%%%%%%%%%%%%%%%%%%%%%%%%%%%%%%
%%%%%%%%%%%%%%%%%%%%%%%%%%%%%%%%%%%%%%%%%%%%%%%%%%%%%%%%%%%%%%%%%%%
% \subsection{MC-O-PI with greedy updates}
\subsection{Guaranteeing convergence with non-uniform updates}
\label{sec: normalisation trick}

The practical takeaway from this chapter is simple: if MC-O-PI seems trapped in a suboptimal limit set (on a boundary or in a cycle), 
then using initial-visit updates
with uniform sampling distributions over the actions in each state
guarantees that the algorithm eventually converges to the optimal policy and its action values
while avoiding slow sliding modes.


% \textcolor{red}{check}
% if compute expectation after sampling initial state-action pair, then can compute update distribution for each state, as just porbability to update it in the episode
% in state st of episode, know pdistribution over at (if use epsilon greedy), then can use that distribution as update probability and cancel it out in learning rate


In order to turn non-uniform update distributions into uniform effective updates, \cite{Tsitsiklis2002OnIteration} suggests making the learning rates inversely proportional to the update distribution.
In a model-free setting, this strategy only applies if the update distributions are generated with the initial-visit method.
This is because the first-visit method causes the update distributions to depend on the transition dynamics induced by the currently greedy policy on the MDP.
Since the greedy policy changes at every step and the structure of the MDP is unknown,
we must assume that first-visit update distributions change at every step.
Cancelling out these changing, non-uniform distributions thus requires knowing the induced transition dynamics, which are unpredictable if the MDP is unknown.

Nonetheless, consider non-uniform, initial-visit update distributions $(\mu_k)_{k \ge 0}$, which essentially equal the sampling distributions $(\sigma_k)_{k \ge 0}$.
For any scalar sequence $(\lambda_k)_{k \ge 0}$ that satisfies the Robbins-Monro conditions,
if the $k^\nth$ episode starts in $(s_0, a_0)$, then setting
\[
    \alpha_k = \frac{\lambda_k}{\sigma_k(s_0, a_0)} , 
\]
guarantees that the effective learning rates in MC-O-PI's difference inclusion are uniform because $\alpha_k \, \mu_k(s_0, a_0) = \lambda_k$.
However, this approach is impractical because it requires knowing the sampling probability for each state-action pair at every step,
or, if these probabilities are constant over time, tracking the number of updates of each state-action pair \citep[p.~66]{Tsitsiklis2002OnIteration}.


% Tsitsiklis notes that MC-O-PI's asymptotic behaviour depends on the asymptotic behaviour of the effective learning rates $(\lrate_k \update_k)$, not on either component alone.
% For instance, if the update distributions are non-uniform but constant, say $\update_k = \update$ for all $k$, and the learning rates take the form $\lrate_k(s,a) = 1/n_k(s,a)$, where $n_k(s,a)$ denotes the number of updates of $(s,a)$ in the first $k$ steps,
% then $\lrate_k \update_k$ converges to $1/k$ because for each state-action pair $n_k(s,a)$ tends to $k \update(s,a)$. As a result, updates are eventually uniform. So \autoref{thm: convergence mcopi uniform} still applies.


In contrast, we showed that global convergence only requires uniform updates over all actions in each state.
Using the decomposition
\[
\sigma_k(s,a) = \sigma^{\sset}_k(s) \, \sigma^{\aset}_k(a \mid s)
\]
where $\sigma^{\sset}$ represents the sampling distribution over the state space,
and $\sigma^{\aset}$ is the distribution over the action space of each state,
it is thus sufficient to choose learning rates as
\[
    \alpha_k = \frac{\lambda_k}{|\mathcal A(s_0)| \ \sigma^{\aset}_k(a_0 \mid s_0)}.
\]
This is more practical because it only assumes knowing the sampling probability of the initial action $a_0$ relative to the other actions of the initial state $s_0$, rather than all other state-action pairs.
We can thus recover global convergence 
% by turning non-uniform update distributions 
% into uniform effective updates 
even if the initial state is sampled from an unknown distribution, provided the distribution of the initial action in that state is known.
The additional term $|\mathcal A(s_0)|$ cancels out $\sigma_k^\aset$ when sampling is already uniform, and it ensures that states do not automatically receive larger updates simply for having more actions.

% This occurs in practice, for instance, when selecting the initial action using an $\varepsilon$-greedy strategy.


% tsitsiklis only works for IV
% our approach works fro FV as well


A concrete example of a non-uniform update distribution arises when biasing updates towards the greedy actions to \textit{exploit} the policy that is currently believed to be optimal.
However, by unevenly updating actions, this approach causes non-greedy actions to evolve very slowly and potentially leads to the sliding modes depicted in \autoref{fig: iv versus fv simulations}.
Therefore, if we additionally scale down the learning rates for the greedy actions in proportion to their increased update probability (so that the product $\alpha_k(s,a)\,\mu_k(s,a)$ does not depend on $a$),
then updates are effectively uniform over the actions of each state, and global convergence is guaranteed.
Notably, this adjustment does not violate the comparability property (\autoref{asp: additional conditions on learning rate}) required by \autoref{thm: convergence mcopi uniform}. Although learning rates may fail global comparability across the entire process, the proof only uses the comparability property between transitions, which matches the intervals in which the greedy actions remain the same.
An additional benefit of making the learning rates inversely proportional to the update distribution is that this reduces the variance of the estimate $q_k$ for greedy actions, which means the solution remains closer to the well-behaved mean-field solution along that dimension.
On the other hand, the larger learning rates for non-greedy actions compensate for their infrequent updates but at the cost of slightly increased variance.
Ultimately, this learning rate adjustment 
% removes the convergence difficulties seen with purely greedy updates (such as slow learning for non-greedy actions, or sliding modes), keeps all trajectories well-behaved, and 
ensures that all action values converge at similar rates along approximately straight trajectories towards the optimum.


% This strategy of scaling the learning rates of a state-action pair inversely with the frequency of updates of that state-action pair
% is similar to the update rule advocated by Sutton and Barto: $\alpha_k(s,a) = 1/n_k(s,a)$, where $n_k(s,a)$ counts the number of times that pair was updated dugin the first $k$ steps.
% This approach automatically results in larger step-sizes for state-action pairs that are updated infrequently, and smaller step-sizes for those that are updated more frequently.
% In the context of non-uniform or greedy-biased updates, this means that the greedy actions, which are selected more often for updates, have their learning rates scaled down because they are temporarily updated more often. Conversely, non-greedy actions, which are temporarily updated less frequently, naturally receive a larger learning rate, so that the (rarer) updates they do receive have a meaningful impact. 
% This ensures that every state-action pair—regardless of how frequently it is updated—receives a fair share in shaping the value estimates, which is precisely what is required for monotonic policy improvement and reliable convergence.
% We suggest a more practical way to essentially do the same thing.


% eg look at the single state but discount => continuiation so that have stochastic returns

% no need to mention
% could maybe do this with first visit + GLIE?
% with GLIE are actually evaluating a stochastic policy (just becomes more and more greedy)
% as epsilon decreases, the policy becomes more and more geredy, so the qpis converge to correct 
% but the update distribution becomes on-hot so issue 
% % 
% still get distribution becoming uniform
% but dthe proof does not work because of dependence of updates

% %%%%%%%%%%%%%%%%%%%%%%%%%%%%%%%%%%%%%%%%%%%%%%%%%%%%%%%%%%%%%%%%%%%
% %%%%%%%%%%%%%%%%%%%%%%%%%%%%%%%%%%%%%%%%%%%%%%%%%%%%%%%%%%%%%%%%%%%
% \subsection{New intuition for Tsitskliks' proof}
% \label{sec: discussion tsitsiklis link}


% \citet{Tsitsiklis2002OnIteration}, \citet{Chen2018OnProblem} and \citet{Liu2021OnStarts} showed that MC-O-PI converges to $\qopt$ when all state-action pairs are updated with equal frequency.
% \autoref{thm: convergence uniform} strictly extends these results by allowing actions of different states to be updated with different frequencies.
% % The proof is also unaffected by the discounting factor as it does not rely on the contraction of Bellman operators.


% Applying the method of \autoref{sec: results} when update distributions are uniform over all state-action pairs 
% reveals a Lyapunov interpretation of Tsitsiklis' framework.
% Namely, consider the potential function that measures the one-sided Bellman residual
% \begin{align}
%     \tilde V(q) = \frac{1}{2} \|(q-\B q)_+\|_2^2 .
% \end{align}
% Its gradient is
% \begin{align}
%     \nabla \tilde V (q) = (I - \discount P_\pi )^\top (q-\B q)_+
% \end{align}
% where $\pi$ is any deterministic policy that is greedy with respect to $q$.
% Computing the inner product with the mean-field update of MC-O-PI yields
% \begin{align*}
%     \langle \nabla \tilde V(q) , q_\pi - q \rangle
%     &= \langle (I - \discount P_\pi )^\top (q-\B q)_+ , (I - \discount P_\pi )^{-1} (\B q - q) \rangle
%     \\
%     &= \langle (q-\B q)_+ , (\B q - q) \rangle
%     \\
%     &= - \|(q - \B q)_+\|_2^2
%     \\
%     &= - 2 \tilde V(q)
% \end{align*}
% Thus, when updates are uniform over all state-action pairs, the mean-field update $(q_\pi - q)$ is a strict descent direction for the potential $\tilde V$.
% This geometric insight corresponds to the decrease in inequality \eqref{eq: q-Bq decreases} that is essential to Tsitsiklis' proof.
% % That proof thus exploits the tension between the fact that every limit point is upper and lower bounded by the best and worst action values of the corresponding greedy policies (\autoref{thm: bounded limit sets})
% % and that such point has a strictly positive one-sided Bellman residual (\autoref{thm: bellman residual has positive component}).
% The proof in \autoref{sec: results} builds on this insight. 
% Unlike Tsitsiklis' method, our proof does not rely on the commutativity of the transition matrix $P_\pi$ and the update distribution $\mu$. The trade-off is that we must instead constrain the variation of the update distribution to ensure that the local descent property is preserved.
% Yet, as discussed earlier, we expect that this conditions can be relaxed.

% We note that Tsitsiklis extends his proof to optimistic policy iteration with temporal difference returns, i.e. where the target $q_\pi$ is replaced by a geometrically weighted sum of multi-step returns (TD($\lambda$)).
% This introduces an additional error term in the descent inequality \eqref{eq: q-Bq decreases}. Given the structural similarities, we conjecture that the framework of \autoref{sec: results} can be similarly extended. Future work could apply our method to TD-O-PI, where an analogous summable error term would likely appear in the critical inequality \eqref{eq: V goes to infinity}, potentially establishing convergence for TD control under our broader class of update distributions.


%%%%%%%%%%%%%%%%%%%%%%%%%%%%%%%%%%%%%%%%%



% For instance, the initial state of each episode can be sampled non-uniformly. As long as the initial action in that state is chosen uniformly, MC-O-PI will still converge to $\qopt$ when using initial-visit updates.

% System \eqref{eq: update unbiased over policies} computes a moving average of the action-values of the policies that it encounters.
% Thus if it encounters a policy $\pol$ a finite number of times, the weight associated with $q_\pol$ monotonically decreases to zero after the last encounter of $\pol$. % call it transient policy?
% In other words, the long-term behaviour of system \eqref{eq: update unbiased over policies} only depends on the policies that it encounters infinitely often.
% Note that, since the number of policies is finite, at least one policy is visited infinitely often.
% \autoref{asp: unbiased over policies} enables us to strengthen this conclusion and show that a single policy is visited infinitely often.


% %%%%%%%%%%%%%%%%%%%%%%%%%%%%%%%%%%%%%%%%%
% \section{Convergence of bootstrapping with uniform bias}

% is not the same as convergence of TD($\lambda$) with uniform bias

% relate to result of \cite{Winnicki2023OnEvaluation}
% Winnicki = no condition on bias but need lookahead (infeasible in practice)
% Us = no need to look ahead but condition on bias


% \vspace{-0.3cm}

%%%%%%%%%%%%%%%%%%%%%%%%%%%%%%%%%%%%%%%%%
%%%%%%%%%%%%%%%%%%%%%%%%%%%%%%%%%%%%%%%%%
\section{Conclusion}

% \vspace{-0.3cm}

% 1. Finding = what did we find
% We established that the mean-field dynamics of MC-O-PI drive trajectories toward monotonically improving policies when updates are uniform over all actions in each state. 
% We then proved that the stochastic perturbations cannot systematically counteract this improvement,
% and consequently that MC-O-PI converges to the optimal action values when updates are uniform.

This chapter established that MC-O-PI converges to the optimal action values when updates are uniform over all actions in each state because solutions transition on average to monotonically improving policies. 
% Consequence of the result
By allowing states to be updated at different frequencies, this result extends existing convergence guarantees to larger-scale problems where uniformly sampling the joint state-action space is intractable.

% Unlike previous proofs that required uniform updates across all state-action pairs, our analysis allows different states to be updated at different frequencies.
% This relaxation makes the algorithm applicable to large-scale problems where uniform sampling over the joint state-action space is computationally intractable.
% We now know that global convergence remains guaranteed as long as it is possible to sample the action space of each state uniformly.
% Instead, learning can prioritise different states at different times by sampling the initial states of episodes accordingly.
% As long as every state is sampled infinitely often and all actions of each state are updated equally frequently, MC-O-PI eventually converges to the optimal policy and action values. 

% \vspace{-0.1cm}


% consequence of the method
We extended the lock-in strategy of the combined stability-ODE method to accommodate the discontinuous dynamics of MC-O-PI.
% \textcolor{red}{Specifically ... does not need to visit strict interior infinitely often, but can visit growing interior infinitely often (then we use seed event to guarantee that enter that growing strict interior)}
This mean-field approach allowed us to analyse the policy sequence $(\pi_k)$ which provided a more intuitive explanation for the convergence mechanism than previous proofs that focused on the estimates $(q_k)$.
More importantly, it opens up new possibilities, beyond contraction-based analysis, to study the convergence of optimistic policy iteration algorithms.

% this proof does not rely on $\gamma<1$ so no need to make difference like previously

% \vspace{-0.1cm}


% perspectives
% The following chapters apply this mean-field approach to study the convergence of MC-O-PI when updates are non-uniform.

% The three main directions for future work are to
% remove the conditions of \autoref{thm: convergence mcopi uniform} that are not necessary for the mean-field improvement property,
% find whether the optimal action values remain globally stable for mean-field solutions of MC-O-PI when updates are non-uniform,
% and apply this new method to study the convergence of optimistic policy iteration algorithms that rely on temporal difference evaluation.


\clearpage